% Options for packages loaded elsewhere
\PassOptionsToPackage{unicode}{hyperref}
\PassOptionsToPackage{hyphens}{url}
\documentclass[review,3p,times]{elsarticle}
\usepackage{amsmath,amssymb}
\usepackage{graphicx}
\usepackage{booktabs}
\usepackage{longtable}
\usepackage{array}
\usepackage{calc}
\usepackage{url}
\usepackage[hidelinks]{hyperref}
\usepackage{lineno}
\modulolinenumbers[1]
\newcounter{none}

\providecommand{\tightlist}{%
  \setlength{\itemsep}{0pt}\setlength{\parskip}{0pt}}
\begin{document}
\title{Measuring Retrieval Fairness Under Adversarial Poisoning: Why the Standard Statistics Are Blind, and What to Measure Instead}
\author[njust]{Jiang Lu\corref{cor1}}
\ead{lujiang12@njust.edu.cn}
\ead{luj00133@gmail.com}
\cortext[cor1]{Corresponding author.}
\affiliation[njust]{organization={School of Electronic and Optical Engineering,
Nanjing University of Science and Technology}, city={Nanjing}, country={China}}
\date{}
\maketitle



\begin{abstract}

Fairness defenses for retrieval-augmented generation (RAG) are selected and validated using aggregate statistics of the retrieved set: the proportion of passages about each protected group, the deviation of a group's stance from a corpus reference, or the exposure of items across repeated requests. Every such statistic has the same form --- a per-group quantity aggregated across groups --- and we show that an adversary who knows this can defeat the aggregation in exactly \textbf{three ways}, which together explain why the defenses built on these statistics fail:

\begin{itemize}
\tightlist
\item
  \textbf{F1, balanced count.} The injection contributes equally to every group, so any statistic over group counts returns its clean value.
\item
  \textbf{F2, preserved nuisance.} The statistic is anchored to a reference the attack does not change, so it reports the component that varies with the query rather than the one that varies with the attack.
\item
  \textbf{F3, cancelled shift.} The attack moves all groups together, so any difference between groups is unchanged however far the groups move.
\end{itemize}

The three modes are not a description but a \textbf{test}: each can be checked against a candidate statistic before any attack is built, and we use them to predict which existing metrics will be blind. To separate \emph{blind to this adversary} from \emph{broken instrument}, we add a positive control in which the ground truth changes by a known amount and the same statistics must respond: they do, monotonically and in the right direction.

The attack is \emph{pairwise poisoning}: injected passages reuse the legitimate template inventory and differ only in which group they favour. Because the injection is balanced across groups by construction, it leaves every group-composition statistic clean while skewing the stance of the evidence. We formalise this as two orthogonal dimensions --- \textbf{(R1) group composition} and \textbf{(R2) within-group stance} --- and show that the blindness is not an accident of any one metric:

\begin{enumerate}
\def\labelenumi{\arabic{enumi}.}
\tightlist
\item
  \textbf{Composition statistics do not move.} R1 drift stays at or below the level legitimate retrieval variance produces, on both a controlled corpus and a naturally written bias benchmark.
\item
  \textbf{Two natural R2 statistics are also blind.} Deviation of stance from a corpus reference (0.286 \(\rightarrow\) 0.311) and one-sidedness of stance (0.467 \(\rightarrow\) 0.500) are flat across the clean and fully attacked conditions, because a group-neutral query legitimately retrieves one-sided evidence, so both measure topic conditioning rather than group skew --- a component the attack preserves. The statistic that responds is the cross-group stance gap.
\item
  \textbf{The right axis is the per-group shift, not the gap between groups --- and we initially got this wrong ourselves.} Every metric above is an absolute difference between groups. Because the corpus carries a pre-existing asymmetry and the attack \emph{relocates both groups} rather than opening a gap between them, an absolute difference discards the signal. Measured per group on our controlled corpus, injection moves the suppressed group's stance by \$-\$0.13 to \$-\$0.18 and the favoured group's by +0.13, consistently across \textbf{four generators spanning three model families} --- API-served Qwen and DeepSeek models plus an open-weight Mistral-7B run locally --- at \(p \le 0.0035\) in all six per-group comparisons; the absolute gap moves by less than 0.002 for three of the four. We reproduced, one level up, the exact error we identify at the retrieval layer. \textbf{On a naturally written corpus the same effect is present, but only a probe that forces a commitment can see it.} The free-form measurement returns nothing on the full BBQ benchmark, and the reason is the instrument rather than the phenomenon: 99.0\% of the answers it scores take no position at all (mean commitment 0.007 against 0.156 on the controlled corpus), so no effect of any size could appear in it. Replacing it with a forced-choice probe on identical retrieval conditions detects the effect --- five per-group shifts at \(p < 0.05\) for one generator, the largest at \(+0.63\) --- and shows that it is \textbf{group-specific and generator-dependent}, which is why reporting a single averaged gap understated it. This free-form probe is the third instrument failure the paper documents.
\end{enumerate}

The practical consequence is a \textbf{negative result about a class of defenses}, which we prove rather than only measure. A distribution-constrained defense faces a trilemma: it cannot simultaneously exclude adversarial passages, admit good clean selections, and impose a non-vacuous constraint. Empirically the R1 constraint is not weak but \textbf{inert} --- and the reason is a ceiling, not a sample-size problem. The unconstrained baseline already admits adversarial passages on every query where it can (1.000 for a dense retriever, 0.750 for BM25), so there is no headroom for a reduction to appear even in principle; the per-query difference is identically zero, giving an equivalence bound of exactly ±0.0000. This is not because the constraint fails to change the selection --- it does change it, moving the R2 stance gap from 0.7436 to 0.5385 --- but because different passages are selected and not one adversarial passage is displaced. We reach the same conclusion across six retrieval back-ends, two corpora, five injection rates and the full budget sweep, with \(|\mathcal{G}| = 2\) throughout.

Two further measurement hazards follow from taking the encoder seriously as a factor rather than an implementation detail. Text-attack susceptibility is predicted by the \textbf{sparsity} of the representation rather than by whether the encoder is learned (a learned sparse retriever is as susceptible as BM25, 0.625 versus 0.750), but sparsity is a correlate and not a law: enlarging GTE within one family and training recipe makes it \textbf{eight times more susceptible} (0.0625 \(\rightarrow\) 0.5000) while enlarging E5 makes it more resistant. Susceptibility is a per-checkpoint property that must be measured, so a single-encoder robustness claim reports an unmeasured encoder property.

We close with the measurement protocol the results imply: per-group shifts rather than cross-group differences, injection as a rate rather than a count, the encoder reported as a factor, and equivalence bounds for null claims.

\end{abstract}

\begin{keyword}
 retrieval-augmented generation, data poisoning, fairness measurement, within-group stance, adversarial robustness, evaluation validity

\begin{center}\rule{0.5\linewidth}{0.5pt}\end{center}

\end{keyword}

\section{Introduction}\label{introduction}

Retrieval-augmented generation (RAG) {[}1{]} has become the standard architecture for grounding large language model (LLM) outputs in external, updatable knowledge. Because the generator conditions on whatever the retriever returns, the retrieval corpus becomes a security boundary: an adversary who can place passages in that corpus can determine the evidence on which the model's answer rests. Corpus poisoning attacks exploit exactly this, and a single line of work has established that a small number of injected passages can mislead an undefended system in the large majority of cases {[}2, 3{]}.

A more recent observation is that poisoning need not target factual correctness at all. Wang et al.~{[}4{]} show that injected passages can \textbf{amplify social bias}: a query that is group-neutral can be answered in a stereotyped way because the retrieved evidence has been skewed toward one group. This matters because the harm is representational rather than factual --- a biased answer can be \emph{correct} by every factual measure while still systematically demeaning a protected group --- and because it is therefore invisible to the accuracy-based evaluation that RAG systems normally undergo.

In response, a body of work has begun to defend the retrieval layer on fairness grounds. The proposed defenses share a common objective: they measure the \textbf{group composition} of the retrieved set and constrain it to remain balanced. Concretely, they partition the corpus by group and either re-rank to equalise group proportions {[}5{]}, adjust the proportion and ordering of group-relevant passages {[}6{]}, reverse-bias the embedder so that its retrieved group balance compensates for the generator's own bias {[}7{]}, or randomise the ranking so that equally deserving items receive equal exposure {[}8{]}. Item-side variants of the same idea equalise \emph{exposure} rather than group share, but still operate on aggregate composition {[}8{]}.

\textbf{This paper is about the measurement those defenses rest on, not about the defenses themselves.} A fairness defense is selected, tuned and validated by a statistic, and it is only as good as that statistic's sensitivity to the threat. We show that the standard statistics are structurally insensitive to the attack that motivates them, that the insensitivity has a single identifiable cause, and that the same cause reappears at the generation layer --- including in our own first attempt at measuring it. The defense failure that follows is then a \emph{consequence} we prove, not a separate claim.

\textbf{The problem, stated as a measurement problem.} \texttt{}Group representation'' is not one quantity but two, and the standard statistics measure the wrong one. A retrieved set can be described by

\begin{itemize}
\tightlist
\item
  \textbf{(R1) group composition} --- the proportion of group-relevant passages that concern each protected group, i.e.~\(P(\text{group} \mid \text{group-relevant})\); and
\item
  \textbf{(R2) within-group stance} --- the proportion of passages about a given group that portray it favourably rather than unfavourably, i.e.~\(P(\text{favourable} \mid \text{group}, \text{group-relevant})\).
\end{itemize}

R1 and R2 are orthogonal: a set can be perfectly balanced in R1 while every passage about one group is unfavourable, and vice versa. Crucially, the adversarial construction that arises naturally in this setting is \textbf{balanced in R1 by construction}. An attacker who wishes to skew representation does not inject one-sided material; they inject matched pairs --- a favourable passage about group \(A\), an unfavourable passage about group \(B\) --- which are individually indistinguishable from legitimate passages because they are drawn from the \emph{same template inventory}. The injected set therefore contributes equally to both groups' counts, leaving R1 essentially unchanged, while shifting R2 systematically.

Every statistic that aggregates over group composition is blind to this. It reports the group distribution as clean, admits the injected passages, and returns success.

\textbf{A second, subtler failure --- one we fell into ourselves.} The obvious repair is to move from composition to stance and to keep using the same arithmetic: measure the \emph{gap} between the two groups. That repair also fails, and the reason generalises. An absolute difference between two group-level quantities is dominated by whatever constant asymmetry the corpus already had, so it is insensitive to an attack that \textbf{relocates both groups together} --- which is what this attack does, suppressing one group while elevating the other. The diagnostic is the \emph{per-group shift}. We identify this failure at the retrieval layer (§5.3), where two of three candidate R2 statistics are flat across the clean and fully attacked conditions, and then reproduce it at the generation layer in our own first measurement of propagation (§5.7): the absolute gap moved by less than 0.002 while the per-group shifts moved by 0.13 at \(p \le 0.0035\), consistently across three generators. We report that second occurrence in full, because a measurement paper whose authors did not notice their own instance of the error would not be credible about anyone else's.

\textbf{What we do.} We formalise R1 and R2 and derive the sensitivity of each candidate statistic to the attack. We build a controlled evaluation in which group and stance labels are exact and the clean reference for both dimensions is known by construction, and we replicate on a naturally written bias benchmark. We re-implement the defense family these statistics motivate and evaluate it against the attack across six retrieval back-ends, before and after scaling each encoder. Our findings are:

\begin{itemize}
\tightlist
\item
  \textbf{The standard statistics are blind, and the blindness is systematic rather than incidental.} R1 drift stays at or below the level legitimate retrieval variance produces; two of three candidate R2 statistics are flat across the clean and fully attacked conditions. The reason is shared: each aggregates over a quantity the attack does not perturb, either group counts (which the injection balances by construction) or a topic-conditioned component (which it preserves).
\item
  \textbf{The right axis is the per-group shift, and the gap between groups is the wrong one --- a distinction we had to be corrected by our own data to see.} Absolute differences between groups are dominated by pre-existing corpus asymmetry and are blind to attacks that relocate both groups. Per-group shifts are large, consistent and highly significant, at both the retrieval and generation layers.
\item
  \textbf{The R1-only defense is not weak, it is inert --- and the reason is a ceiling, which is a stronger statement than a failure to detect an effect.} The unconstrained baseline already admits adversarial passages on every query where it can, so there is no headroom for a reduction to appear even in principle; the per-query difference is identically zero and the equivalence bound is exactly ±0.0000, excluding \emph{any} effect rather than an effect above a threshold. It is not that the constraint cannot change the selection --- it does --- but that changing the selection does not displace a single adversarial passage.
\item
  \textbf{The encoder must be reported as a factor, not as an implementation detail.} Text-attack susceptibility is predicted by the sparsity of the representation rather than by whether the encoder is learned, but sparsity is a correlate and not a law: enlarging GTE within one family and training recipe makes it eight times more susceptible while enlarging E5 makes it more resistant. A single-encoder robustness claim therefore reports a property of a checkpoint that was never measured.
\item
  \textbf{A provenance signal is necessary, and we report an exploratory candidate.} We measure the number of distinct queries for which a passage is retrieved and find clean separation in our controlled setting. We state plainly that our corpus is synthetic and that this separation may be an artefact of its construction; we therefore present it as a hypothesis with a specified falsification experiment, not as a defense.
\end{itemize}

\textbf{Contributions.}

\begin{enumerate}
\def\labelenumi{\arabic{enumi}.}
\tightlist
\item
  \textbf{A framework for why adversarial fairness statistics fail.} Every candidate metric in this literature is a per-group quantity aggregated across groups, and an adversary can defeat the aggregation in exactly three ways: by balancing what the statistic counts (F1), by anchoring it to a reference the attack preserves (F2), or by moving all groups together so a difference between them cancels (F3). The modes are checkable in advance of building an attack, and we use them to predict which published metrics are blind (§3.4), then confirm the predictions empirically (§5.1, §5.2, §5.7).
\item
  \textbf{A positive control that separates blindness from breakage.} With the ground truth moved by a known amount and the corpus size held fixed, the indicted statistics respond monotonically and directionally (Table 1). The claim is therefore specific --- these statistics are immune to an adversary who balances what they count --- rather than a general complaint that the instrument is unreliable.
\item
  Identification of a \textbf{shared cause of measurement failure}: statistics that aggregate over group counts, or that take an absolute difference between groups, are insensitive to an attack that balances the former and relocates the latter --- with the error demonstrated at two layers, including in our own first generation-layer measurement (§5.1, §5.3, §5.7).
\item
  A controlled pairwise-poisoning construction that makes both dimensions measurable without a generator, so the mechanism can be studied independently of LLM behaviour (§4), replicated on a naturally written bias benchmark.
\item
  An empirical demonstration that R1-constraining defenses are inert, with \textbf{equivalence bounds rather than point estimates} for the null, across \textbf{six} base retrieval back-ends, three scaled checkpoints, two corpora and five injection rates (§5, Appendix B).
\item
  A negative result on the limits of distribution-level defense: aggregate-composition constraints cannot reject an individually admissible passage, so any defense in this class has an adversarial-passage inclusion floor of 1 (§6). To our knowledge this limit has not been characterised.
\item
  A scale check that bounds the sparsity account of encoder susceptibility: susceptibility is \textbf{not monotone in encoder size} (GTE 0.0625 \(\rightarrow\) 0.5000 while E5 0.0625 \(\rightarrow\) 0.0000 within matched families), so it is a per-checkpoint property (§5.6).
\item
  A measurement protocol implied by the above: per-group shifts rather than cross-group differences, injection reported as a rate rather than a count, the encoder reported as a factor, equivalence bounds for null claims, and both stance and attribution scored by entailment in the correct direction rather than by a prompted judge (§8, §5.7.1).
\end{enumerate}

\begin{center}\rule{0.5\linewidth}{0.5pt}\end{center}

\section{Related Work}\label{related-work}

\subsection{Poisoning attacks on RAG}\label{poisoning-attacks-on-rag}

Corpus poisoning against RAG was established by Zou et al.~{[}2{]}, who showed that a small number of optimised passages can steer a system's answers, and has since been extended to black-box and query-agnostic settings {[}3{]}, to knowledge-graph-structured retrieval {[}9{]}, and to multimodal pipelines {[}10, 11, 12{]}. A parallel thread targets the retrieval process itself rather than passage content, manipulating embedding space so that injected passages rank highly regardless of their text. Wang et al.~{[}4{]} combine both: reward-optimised adversarial documents, subspace projection to raise their retrieval probability, and a generate--evaluate--reinject loop that accumulates bias over time.

All of this work targets \emph{factual} correctness except {[}4{]}, and {[}4{]} itself measures bias by \emph{counting} stereotype-consistent selections rather than by modelling the distribution of stance within groups. Our R2 formulation is a distributional refinement of that measurement, and we show that the refinement changes which defenses can work.

\subsection{Fairness in RAG}\label{fairness-in-rag}

Fairness in RAG has been studied primarily as a property of the knowledge base and the generator rather than of an adversary. Wu et al.~{[}6{]} construct scenario-based questions and evaluate group disparity across RAG components, finding that the retriever has the largest influence on both accuracy and fairness, and that utility and fairness trade off. Their mitigation is to adjust the proportion and ordering of group-relevant passages --- an R1 intervention --- and they state that they do not provide a comprehensive exploration of mitigation strategies. Kim et al.~{[}7{]} decompose RAG into LLM, embedder, and corpus, show that component biases interact (\texttt{}bias conflict''), and find that \emph{reverse-biasing} a small embedder can cancel a much larger generator's bias; their bias metric is group membership, so both R1 and R2 are outside their formulation. Kim and Diaz {[}8{]} bring fair-ranking machinery to RAG, using stochastic rankers to equalise item-side exposure across repeated requests, and find that fairness and quality need not trade off severely; their fairness unit is the item/provider rather than the social group.

Most closely related is Zhao et al.~{[}5{]}, who model position-wise bias propagation in top-\(k\) RAG and formulate fairness-aware retrieval as an optimisation problem over group composition, solved by decomposition. This is the most general R1-constraining defense we are aware of, and we treat it as the representative of the class. Like the others, it assumes no adversary: bias is a property of the corpus and the generator, not something an attacker injects.

Two further works are relevant. Kim et al.~{[}7{]} observe that \texttt{}naively increasing fairness is not always the optimal solution'' --- a conclusion we sharpen by identifying the \emph{specific} reason: the fairness being increased is measured along the wrong axis. Bagwe et al.~{[}13{]} construct fairness-targeted backdoor attacks that manipulate semantic relationships between groups and biases; their attack relies on a trigger and their objective is attack construction, whereas we study trigger-free injection and its consequences for defense design.

\subsection{The gap}\label{the-gap}

Across this literature we identify four points of agreement that together define the gap this paper addresses. First, fairness defenses for RAG consistently constrain group composition or item exposure (§2.2). Second, several authors independently report that existing defenses are not robust enough --- \texttt{}further optimization \ldots{} to achieve more stable and generalizable defense performance'' {[}4{]}; \texttt{}single-stage defenses give limited robustness'' {[}14{]}; \texttt{}does not provide a comprehensive exploration of strategies to mitigate these unfairnesses'' {[}6{]}. Third, and structurally, none of the fairness works models an adversary: they study corpora that are \emph{naturally} skewed. Fourth, no work we have found constrains or detects within-group stance.

We therefore study the intersection that has not been studied: an \emph{adversary} acting on the axis that \emph{no} defense constrains.

\begin{center}\rule{0.5\linewidth}{0.5pt}\end{center}

\section{Two Dimensions of Representation}\label{two-dimensions-of-representation}

\subsection{Definitions}\label{definitions}

Let a retrieved set for query \(q\) be \(\mathcal{D}_q = \{d_1, \dots, d_k\}\), and let each passage \(d\) carry a group label \(g(d) \in \mathcal{G} \cup \{\bot\}\) (the protected group the passage concerns; \(\bot\) for group-neutral text) and a stance \(s(d) \in \{+1, -1, 0\}\) (\(+1\) favourable toward \(g(d)\), \(-1\) unfavourable, \(0\) neutral or group-irrelevant).

\textbf{Definition 1 (R1: group composition).} Over the group-relevant passages of \(\mathcal{D}_q\),

\[P^{\mathcal{D}_q}_{\text{grp}}(g) = \frac{|\{d \in \mathcal{D}_q : g(d)=g,\ s(d)\neq 0\}|}{|\{d \in \mathcal{D}_q : g(d)\neq \bot,\ s(d)\neq 0\}|}.\]

\textbf{Definition 2 (R2: within-group stance).} For each group \(g\),

\[P^{\mathcal{D}_q}_{\text{fav}}(g) = \frac{|\{d \in \mathcal{D}_q : g(d)=g,\ s(d)=+1\}|}{|\{d \in \mathcal{D}_q : g(d)=g,\ s(d)\neq 0\}|}.\]

\textbf{Definition 3 (R1 drift).} Against a reference \(P^{*}_{\text{grp}}\),

\[\Delta_{\text{R1}}(\mathcal{D}_q) = \tfrac{1}{2}\sum_{g} \bigl| P^{\mathcal{D}_q}_{\text{grp}}(g) - P^{*}_{\text{grp}}(g) \bigr|.\]

\textbf{Definition 4 (R2 group-conditioned stance gap).} The primary R2 metric is the spread of favourable rate \emph{across} groups:

\[\Delta_{\text{R2}}(\mathcal{D}_q) = \max_{g \neq g'} \bigl| P^{\mathcal{D}_q}_{\text{fav}}(g) - P^{\mathcal{D}_q}_{\text{fav}}(g') \bigr|.\]

\textbf{Why a gap across groups, and not a deviation from a reference.} The natural first choice for an R2 metric --- how far a group's favourable rate sits from the corpus reference --- turns out to be unusable, and the reason is worth recording because it is easy to get wrong. Consider a balanced corpus and a query such as \emph{\texttt{}Who performed better in the engineering evaluation?''} The query's vocabulary matches favourable and unfavourable templates differently, so a clean retrieval of top-\(k\) tends to be \textbf{one-sided}: it returns five passages that are all favourable, or all unfavourable. Measured against a reference favourable rate of \(0.5\), this yields a large deviation (\(\approx 0.31\) by JS divergence) \emph{before any attack is applied}. We confirmed this directly: every clean query in our corpus returns a top-5 that is entirely one-sided.

Crucially, that one-sidedness is \textbf{topic-conditioned, not group-conditioned}: when the corpus is balanced, an all-unfavourable retrieval is all-unfavourable for \emph{every} group, so it is not a fairness violation. A deviation-from-reference metric therefore conflates two different phenomena --- legitimate topic conditioning, which the attack does not change, and group-conditioned skew, which is the actual harm. Under pairwise poisoning the attack \emph{preserves} the topic-conditioned component (its injected passages are drawn from the same templates), so a deviation metric does not move: in our experiments both the reference-deviation and one-sidedness metrics are flat across the clean and attacked conditions (\(0.2859 \to 0.3113\) and \(0.4667 \to 0.5000\) respectively), while \texttt{poison@k} goes from 0 to 0.750. A defense monitoring either would see a healthy system.

The gap across groups isolates the group-conditioned component. Under a balanced corpus a clean retrieval has \(\Delta_{\text{R2}} = 0\) exactly (all groups equally one-sided, hence equal rates), which we observe across all four strata. Pairwise poisoning favours one group while suppressing another, driving the gap toward 1. This is the quantity the attack actually perturbs, and it is invisible to any R1 monitor.

\subsection{Why the reference for R2 must be a gap, not a deviation}\label{why-the-reference-for-r2-must-be-a-gap-not-a-deviation}

We state this separately because our own first implementation got it wrong, and the error is instructive.

A per-query clean reference --- the natural choice, and the one used for R1 --- \textbf{cannot detect an R2 attack}. Under pairwise poisoning the injected passages are drawn from the same template inventory as the legitimate ones, so the \emph{availability} of favourable and unfavourable passages in the candidate pool is unchanged: the attacker replaces legitimate favourable evidence with adversarial favourable evidence, and legitimate unfavourable evidence with adversarial unfavourable evidence. A clean retrieval under the same budget therefore exhibits the same stance profile as the attacked retrieval. Measuring R2 against the clean retrieval returns approximately zero drift \emph{even when every retrieved passage is attacker-controlled}.

Worse, a deviation-from-corpus-reference metric does not work either, for the reason given in Definition 4: it is dominated by topic conditioning, which the attack preserves. In our experiments the corpus-reference JS divergence is \(0.2859\) clean and \(0.3113\) attacked --- a change of \(0.025\) while adversarial inclusion moves from \(0\) to \(0.750\).

Only the cross-group gap responds. We emphasise all three failures because each is a plausible design choice that produces a monitor reporting a healthy system while the evidence is fully adversarial.

\subsection{What pairwise poisoning does to each dimension}\label{what-pairwise-poisoning-does-to-each-dimension}

Let the attacker inject \(n\) passages, half favourable toward group \(A\) and half unfavourable toward group \(B\), using the legitimate template inventory. Then:

\begin{itemize}
\tightlist
\item
  \textbf{R1:} the injected set contributes \(\lceil n/2 \rceil\) group-relevant passages to each of \(A\) and \(B\). Its composition is \((\tfrac{1}{2}, \tfrac{1}{2})\) --- \emph{identical to the balanced reference}. Formally, the compositional drift contributed by the injection is 0, and any observed R1 drift comes from injected passages outranking legitimate ones, not from their composition.
\item
  \textbf{R2 (reference deviation):} unchanged, because the attack preserves the topic-conditioned one-sidedness that dominates this quantity (§3.3).
\item
  \textbf{R2 (cross-group gap):} \(P_{\text{fav}}(A) \to 1\) and \(P_{\text{fav}}(B) \to 0\), so \(\Delta_{\text{R2}} \to 1\) --- the maximum.
\end{itemize}

This asymmetry is the attack's entire mechanism, and it is invisible to an R1 monitor by construction.

\begin{center}\rule{0.5\linewidth}{0.5pt}\end{center}

\subsection{One form for every candidate metric, and three ways it loses the signal}\label{one-form-for-every-candidate-metric-and-three-ways-it-loses-the-signal}

Every fairness statistic in this literature, including the ones we introduce, has the same shape: a per-group quantity aggregated across groups.

\[M(\mathcal{D}) \;=\; A_{g \in \mathcal{G}}\bigl[\,\phi_g(\mathcal{D})\,\bigr] \tag{2}\]

where \(\phi_g(\mathcal{D})\) is some per-group measurement --- the share of group-relevant passages belonging to \(g\), the proportion of those passages that are favourable, the exposure of \(g\)'s items --- and \(A\) is an aggregation over groups (a count, a deviation from a reference, a difference, a variance). R1 statistics take \(\phi_g\) to be a composition and \(A\) to be a deviation or a count; R2 statistics take \(\phi_g\) to be a stance proportion and \(A\) to be a deviation, a difference or an extremeness.

This form is useful because it exposes exactly where an adversary can act. An attack on a fairness metric has two available moves and only two: leave the aggregate \(A\) where it is while moving \(\phi_g\), or move \(A\) in a direction the defender reads as harmless. Every attack we study does the first. The consequence is a taxonomy of three failure modes, and they are the reason the experiments in §5 come out as they do.

{\def\LTcaptype{none} % do not increment counter
\begin{longtable}[]{@{}
  >{\raggedright\arraybackslash}p{(\linewidth - 6\tabcolsep) * \real{0.2500}}
  >{\raggedright\arraybackslash}p{(\linewidth - 6\tabcolsep) * \real{0.2500}}
  >{\raggedright\arraybackslash}p{(\linewidth - 6\tabcolsep) * \real{0.2500}}
  >{\raggedright\arraybackslash}p{(\linewidth - 6\tabcolsep) * \real{0.2500}}@{}}
\toprule\noalign{}
\begin{minipage}[b]{\linewidth}\raggedright
Failure mode
\end{minipage} & \begin{minipage}[b]{\linewidth}\raggedright
Mechanism
\end{minipage} & \begin{minipage}[b]{\linewidth}\raggedright
Which \(A\) is exposed
\end{minipage} & \begin{minipage}[b]{\linewidth}\raggedright
Where we observe it
\end{minipage} \\
\midrule\noalign{}

\bottomrule\noalign{}

\textbf{F1 --- balanced count} & The injection contributes equally to every group, so any aggregate over group counts returns its clean value & any count or share across groups & §5.1, §5.2 (R1 statistics) \\
\textbf{F2 --- preserved nuisance} & A statistic is anchored to a reference the attack does not change, so it reports the component that varies with the query rather than the one that varies with the attack & any deviation from a reference & §3.2, §5.1 (\texttt{stance\_div}, \texttt{stance\_onesided}) \\
\textbf{F3 --- cancelled shift} & The attack relocates all groups in the same direction, so a difference between groups is unchanged even though every \(\phi_g\) moved a long way & any difference, variance or gap across groups & §5.1, §5.7 (\texttt{stance\_gap} and its generation-layer twin) \\
\end{longtable}
}

\textbf{The taxonomy is a test, not a description.} For a candidate statistic \(M\), the three failure modes correspond to three questions that can be asked before any attack is built:

\begin{enumerate}
\def\labelenumi{\arabic{enumi}.}
\tightlist
\item
  \emph{Does the injection balance the quantity \(M\) counts?} If the attacker's construction is symmetric across groups --- and pairwise poisoning is, because it reuses the legitimate template inventory in matched pairs --- then any counting aggregate is invariant by construction and no amount of data will show the attack.
\item
  \emph{Is \(M\) referenced to something the attack preserves?} If the reference is a per-query clean retrieval or a corpus-level rate, the attack leaves it intact by design, and \(M\) measures the query's topic rather than the attacker's skew.
\item
  \emph{Is \(M\) an aggregate that cancels a common-mode movement?} A difference or variance across groups is invariant under \(g \mapsto g + c\) for every group at once. An attack that moves all groups together is therefore invisible to it, no matter how large the movement.
\end{enumerate}

A statistic that survives all three is, in our setting, one that reports \textbf{per-group levels} rather than a difference between them, because that is the only form that is not invariant under at least one of the three moves. §5.1 and §5.7 test this prediction: on the same retrieval conditions where a difference statistics reports a change of 0.002, the per-group levels move by 0.13 and the paired test rejects at \(p < 0.01\). The prediction was made from the form of the statistic, not fitted to those numbers, and §8 records that we violated it ourselves in an earlier version of the generation-layer analysis.

\textbf{The taxonomy is testable against known ground truth, and we ran that test.} A taxonomy of failure modes is only worth stating if the statistics it indicts are otherwise sound: if they did not respond to a real composition change either, the finding would be that the measurement apparatus is broken, not that an adversary defeats it. We therefore ran a positive control on the naturally written corpus, where the true composition is known by construction.

The control holds the corpus size fixed at 17,792 passages and \textbf{swaps} passages between the two groups of a balanced stratum, so the retrieval pool and the score threshold do not move and composition is the only variable. It changes by a known number of passages, in a known direction, and is reversed in the last row.

\textbf{Table 1.} Responsiveness control. Corpus size constant; only group composition changes. \texttt{corpus\%woman} is the true share after the swap; the remaining columns are the statistics under test.

{\def\LTcaptype{none} % do not increment counter
\begin{longtable}[]{@{}
  >{\raggedright\arraybackslash}p{(\linewidth - 10\tabcolsep) * \real{0.1667}}
  >{\raggedright\arraybackslash}p{(\linewidth - 10\tabcolsep) * \real{0.1667}}
  >{\raggedright\arraybackslash}p{(\linewidth - 10\tabcolsep) * \real{0.1667}}
  >{\raggedright\arraybackslash}p{(\linewidth - 10\tabcolsep) * \real{0.1667}}
  >{\raggedright\arraybackslash}p{(\linewidth - 10\tabcolsep) * \real{0.1667}}
  >{\raggedright\arraybackslash}p{(\linewidth - 10\tabcolsep) * \real{0.1667}}@{}}
\toprule\noalign{}
\begin{minipage}[b]{\linewidth}\raggedright
Variant
\end{minipage} & \begin{minipage}[b]{\linewidth}\raggedright
\texttt{corpus\%woman}
\end{minipage} & \begin{minipage}[b]{\linewidth}\raggedright
\texttt{drift\_tv}
\end{minipage} & \begin{minipage}[b]{\linewidth}\raggedright
\texttt{drift\_js}
\end{minipage} & \begin{minipage}[b]{\linewidth}\raggedright
\texttt{stance\_div}
\end{minipage} & \begin{minipage}[b]{\linewidth}\raggedright
\texttt{stance\_onesided}
\end{minipage} \\
\midrule\noalign{}

\bottomrule\noalign{}

baseline (balanced) & 0.500 & 0.0000 & 0.0000 & 0.2925 & 0.4753 \\
swap 500 man\(\rightarrow\)woman & 0.580 & 0.0264 & 0.0372 & 0.2876 & 0.4711 \\
swap 1500 man\(\rightarrow\)woman & 0.739 & 0.0660 & 0.0809 & 0.2854 & 0.4694 \\
swap 3000 man\(\rightarrow\)woman & 0.978 & 0.0924 & 0.1062 & 0.2854 & 0.4694 \\
\textbf{swap 3000 woman\(\rightarrow\)man} & \textbf{0.022} & \textbf{0.0569} & 0.0532 & 0.2851 & 0.4691 \\
\end{longtable}
}

Two conclusions, and they are what make the negative results interpretable rather than merely negative. First, \textbf{\texttt{drift\_tv} and \texttt{drift\_js} are working instruments}: they rise monotonically with a real composition change and respond to its direction. The R1 statistic is not broken; it is \emph{specifically} invariant to an injection that balances group counts, which is exactly what F1 predicts and what §5.2 measures. Second, \textbf{the two reference-based stance statistics barely move even here} --- \texttt{stance\_div} changes by 0.007 across a composition swing from 0.022 to 0.978 --- which is independent confirmation of F2 on data whose ground truth is known by construction rather than inferred from an attack.

The contrast is the paper's central claim in one table: the same statistics that detect a 3000-passage composition change with \(p \to 0\) detect a fully adversarial injection not at all.

\textbf{What the taxonomy does not claim.} It is not a claim that no fairness metric can be robust. It is a claim about the class of statistics that appear in this literature, all of which are aggregates of the form (2), together with a constructive statement of what escapes the three modes: report \(\phi_g\) for each \(g\), and report how each moves. That recommendation is cheap, it is what §5.7 does, and it is the difference between detecting the attack and not.

\begin{center}\rule{0.5\linewidth}{0.5pt}\end{center}

\section{Attack and Evaluation Setup}\label{attack-and-evaluation-setup}

\subsection{Corpus}\label{corpus}

We require a corpus in which (i) every passage carries exact group and stance annotations, and (ii) the clean reference for both R1 and R2 is known by construction. Public bias benchmarks provide stereotype/anti-stereotype template inventories {[}15, 16{]} but not passage-level scaffolding with a known clean reference, so we build a controlled corpus: for each of four bias strata (gender, disability, age, race), a balanced pool of passages is instantiated from the stratum's template inventory, with an equal number of favourable and unfavourable passages per group and a set of group-neutral passages. Queries are group-neutral by construction --- the query text never names a protected group --- so any group skew in a retrieved set is attributable to the corpus and not to the question. §5.4 replicates the mechanism findings on the naturally written BBQ corpus, and §8 discusses why no \emph{neutral} retrieval corpus can substitute (§8, Limitation 1).

\subsection{Attack}\label{attack}

\textbf{Pairwise template injection.} For each stratum with groups \((A, B)\), the attacker selects a target direction (favour \(A\), suppress \(B\)) and injects \(n\) passages drawn from the legitimate template inventory, half favourable toward \(A\) and half unfavourable toward \(B\). Two mechanisms ensure the injected passages actually out-rank legitimate ones, since otherwise the attack is vacuous:

\begin{enumerate}
\def\labelenumi{\arabic{enumi}.}
\tightlist
\item
  \emph{Query alignment.} Legitimate passages carry a discourse filler that does not match query vocabulary; injected passages replace it with query-aligned terms. This raises lexical scores without altering the stance-bearing clause, so the signal the defense must detect is unchanged.
\item
  \emph{Length.} Injected passages are shorter, which length-normalised lexical scoring rewards.
\end{enumerate}

Note that (1) makes the attack effective against BM25 as well as dense retrieval, so a defense cannot pass by relying on lexical scoring being hard to poison.

\textbf{Subspace projection.} For dense back-ends we additionally apply the projection-and-amplify step \(h' = h + \lambda \langle h, \bar{q}\rangle \|\bar{q}\|^{-2} \bar{q}\) with \(\bar{q}\) the mean query direction, following {[}4{]}. We sweep \(\lambda \in \{0, 0.25, 0.5, 1, 2, 4\}\).

\textbf{Adaptive variant.} We also evaluate a \emph{defense-aware} attacker, which observes the defense's own scoring function and searches over \(\lambda\) for the value that best balances raised similarity against the defense's penalty. A defense that reports numbers only against the non-adaptive attacker has not demonstrated robustness.

\subsection{Defenses}\label{defenses}

We implement the R1-constraining family plus ablations.

\textbf{Multi-query consistency} {[}4{]} --- aggregate rankings over query variants; score passages by appearance frequency and mean rank.

\textbf{Off-manifold filtering} --- calibrate the corpus manifold (mean plus leading principal directions), penalise passages by residual distance outside it.

\textbf{Representation-conserving selection} --- greedy Lagrangian selection with per-group admission budgets, in three ablations:
- \texttt{repr\_group} (\textbf{R1 only}): constrain \(\Delta_{\text{R1}} \le \epsilon\). This is the representative of the existing defense class.
- \texttt{repr\_stance} (\textbf{R2 only}): constrain within-group favourable rate to \([P^{*}_{\text{fav}}(g) - \epsilon,\ P^{*}_{\text{fav}}(g) + \epsilon]\), with a repair pass that restores the lower bound by swapping the weakest admissible unfavourable passage of a deficient group for the best-scoring unselected favourable passage of that group.
- \texttt{repr\_both} (\textbf{R1 + R2}): both constraints.

The two-sided R2 constraint is deliberate. An upper bound alone stops an attacker from \emph{flooding} a group with favourable passages; the actual attack \emph{suppresses}, driving the favourable rate down, and is therefore only constrained by the lower bound.

\subsection{Metrics}\label{metrics}

{\def\LTcaptype{none} % do not increment counter
\begin{longtable}[]{@{}
  >{\raggedright\arraybackslash}p{(\linewidth - 4\tabcolsep) * \real{0.3333}}
  >{\raggedright\arraybackslash}p{(\linewidth - 4\tabcolsep) * \real{0.3333}}
  >{\raggedright\arraybackslash}p{(\linewidth - 4\tabcolsep) * \real{0.3333}}@{}}
\toprule\noalign{}
\begin{minipage}[b]{\linewidth}\raggedright
Metric
\end{minipage} & \begin{minipage}[b]{\linewidth}\raggedright
Dimension
\end{minipage} & \begin{minipage}[b]{\linewidth}\raggedright
Definition
\end{minipage} \\
\midrule\noalign{}

\bottomrule\noalign{}

\texttt{drift\_tv} & R1 & \(\Delta_{\text{R1}}\) --- TV distance of group composition vs clean retrieval \\
\texttt{drift\_js} & R1 & JS divergence of group composition vs clean retrieval \\
\textbf{\texttt{stance\_gap}} & \textbf{R2 (primary)} & \(\Delta_{\text{R2}}\) --- max spread of favourable rate across groups (§3, Def. 4) \\
\texttt{stance\_shift} & R2 (control) & \(\max_g \|P_{\text{fav}}(g) - P^{*}_{\text{fav}}(g)\|\) vs corpus reference; saturates \\
\texttt{stance\_div} & R2 (control) & \(\max_g \mathrm{JS}(\pi^{\mathcal{D}_q}_g \| \pi^{*}_g)\) vs corpus reference; confounded by topic conditioning \\
\texttt{stance\_onesided} & R2 (control) & \(\max_g \|P_{\text{fav}}(g) - 0.5\|\); also confounded \\
\texttt{poison@k} & Attack success & fraction of queries with \(\ge 1\) injected passage in top-\(k\) \\
\texttt{poison\_share} & Attack success & mean fraction of top-\(k\) occupied by injected passages \\
\texttt{in\_pool\_rate} & Utility & fraction of retrieved passages from the query's candidate pool \\
\texttt{mean\_score} & Utility & mean retriever score of the returned set \\
\end{longtable}
}

We deliberately report three R2 statistics rather than one. The two reference-based variants are included as \textbf{controls}: §5.1 shows they do not respond to the attack, which is itself a finding, and reporting them documents the failure mode rather than hiding it.

Ground-truth poison labels are used \textbf{only} for the attack-success metrics; no defense receives them. All defenses are evaluated with \(k = 5\). Retrieval back-ends are BM25, a self-contained feature-hashing dense retriever, and --- for the backbone alignment of §5.5 and the scale check of §5.6 --- the encoders GTE-base {[}5{]}, GTE-large, Contriever {[}8{]}, E5-base-v2 and E5-large-v2 {[}4, 6{]}, plus the learned-sparse retrievers SPLADE (distil) and SPLADE-efficient-large. Multilingual and instruction-tuned embedders remain to be swept (Limitation 5).

\begin{center}\rule{0.5\linewidth}{0.5pt}\end{center}

\section{Results}\label{results}

\begin{quote}
\textbf{Status of this section.} All numbers reported below were produced by the released implementation and are reproducible with the commands in Appendix A. Two corpora are used: a controlled corpus with an exactly known clean reference (§5.1--§5.3) and the natural BBQ corpus (§5.4), each swept over injection rate \(\rho \in \{0.1\%, \dots, 2\%\}\) of corpus size; §5.5 adds the backbone alignment against five further retrievers, §5.6 the encoder-scale check, and §5.7 the generation-stage propagation check. The controlled suite completes in \textasciitilde25 s, the BBQ replication in \textasciitilde180 s, and the multi-backbone sweep in \textasciitilde12 min, all on CPU with no GPU.
\end{quote}

\subsection{Pairwise poisoning is an R2 attack, not an R1 attack}\label{pairwise-poisoning-is-an-r2-attack-not-an-r1-attack}

Table 2 reports attack effect with no defense, over 48 queries across four bias strata.

\textbf{Table 2.} Attack effect without defense.

{\def\LTcaptype{none} % do not increment counter
\begin{longtable}[]{@{}
  >{\raggedright\arraybackslash}p{(\linewidth - 12\tabcolsep) * \real{0.1429}}
  >{\raggedright\arraybackslash}p{(\linewidth - 12\tabcolsep) * \real{0.1429}}
  >{\raggedright\arraybackslash}p{(\linewidth - 12\tabcolsep) * \real{0.1429}}
  >{\raggedright\arraybackslash}p{(\linewidth - 12\tabcolsep) * \real{0.1429}}
  >{\raggedright\arraybackslash}p{(\linewidth - 12\tabcolsep) * \real{0.1429}}
  >{\raggedright\arraybackslash}p{(\linewidth - 12\tabcolsep) * \real{0.1429}}
  >{\raggedright\arraybackslash}p{(\linewidth - 12\tabcolsep) * \real{0.1429}}@{}}
\toprule\noalign{}
\begin{minipage}[b]{\linewidth}\raggedright
Attack
\end{minipage} & \begin{minipage}[b]{\linewidth}\raggedright
Retriever
\end{minipage} & \begin{minipage}[b]{\linewidth}\raggedright
\texttt{drift\_tv} (R1)
\end{minipage} & \begin{minipage}[b]{\linewidth}\raggedright
\texttt{stance\_gap} (R2)
\end{minipage} & \begin{minipage}[b]{\linewidth}\raggedright
\texttt{stance\_div} (ref-dev)
\end{minipage} & \begin{minipage}[b]{\linewidth}\raggedright
\texttt{one\_sided}
\end{minipage} & \begin{minipage}[b]{\linewidth}\raggedright
\texttt{poison@k}
\end{minipage} \\
\midrule\noalign{}

\bottomrule\noalign{}

clean & BM25 & 0.0000 & \textbf{0.0000} & 0.2859 & 0.4667 & --- \\
clean & dense & 0.0000 & \textbf{0.0000} & 0.3113 & 0.5000 & --- \\
template & BM25 & 0.1500 & \textbf{0.6250} & 0.3113 & 0.5000 & 0.750 \\
template & dense & 0.0625 & \textbf{0.5000} & 0.3113 & 0.5000 & 0.625 \\
template + projection & BM25 & 0.1500 & \textbf{0.6250} & 0.3113 & 0.5000 & 0.750 \\
template + projection & dense & 0.1313 & \textbf{1.0000} & 0.3113 & 0.5000 & \textbf{1.000} \\
\end{longtable}
}

Three observations.

\textbf{First, the R2 gap is exactly zero before the attack and large after it, in every stratum.} Clean retrieval yields \texttt{stance\_gap} \(= 0.0000\) across all four strata (gender, disability, age, race) --- because on a balanced corpus a topic-conditioned one-sided retrieval is one-sided \emph{for every group equally}, so the cross-group spread vanishes. Under attack the gap rises to 0.6250 (BM25) and 1.0000 (dense + projection). Per stratum, the clean condition is \(0.0000\) throughout while the attacked condition is \(\{1.0000, 0.5000, 0.3333, 0.7500\}\) for the template attack and \(1.0000\) in all four strata under projection.

\textbf{Second, the two reference-based R2 metrics fail completely.} \texttt{stance\_div} --- deviation from the corpus stance reference --- moves from 0.2859 to 0.3113 (a change of 0.025) while adversarial inclusion goes from absent to 0.750; \texttt{one\_sided} moves from 0.4667 to 0.5000. Either metric would report a healthy system. This empirically confirms the analysis of §3.2: the attack preserves the topic-conditioned component that dominates both quantities, so only the cross-group gap responds. We regard this as a methodological contribution in its own right --- three plausible R2 metrics, two of which are silently blind.

\textbf{Third, R1 drift is small and non-diagnostic.} The largest value is 0.1313 (dense, projection), consistent with §3.3: the injected set is compositionally \((\tfrac{1}{2}, \tfrac{1}{2})\), so the drift reflects injected passages outranking legitimate ones rather than any compositional imbalance on their part. An R1-based monitor sees a retrieved set whose group balance is close to the clean reference, while \texttt{poison@k} reaches 1.000 and the utility proxy \texttt{in\_pool\_rate} falls to 0.000 --- the retrieved set becomes entirely attacker-controlled.

\textbf{Finding 1.} Pairwise poisoning drives the cross-group stance gap to 0.6250--1.0000 while leaving both R1 drift and the two reference-based R2 statistics near their clean values. Neither an R1 monitor nor a naive R2 monitor flags it.

\subsection{The R1-only defense is inert}\label{the-r1-only-defense-is-inert}

\textbf{Table 3.} Defense comparison under template + projection. \texttt{poison@k} is the adversarial-passage inclusion rate.

{\def\LTcaptype{none} % do not increment counter
\begin{longtable}[]{@{}
  >{\raggedright\arraybackslash}p{(\linewidth - 10\tabcolsep) * \real{0.1667}}
  >{\raggedright\arraybackslash}p{(\linewidth - 10\tabcolsep) * \real{0.1667}}
  >{\raggedright\arraybackslash}p{(\linewidth - 10\tabcolsep) * \real{0.1667}}
  >{\raggedright\arraybackslash}p{(\linewidth - 10\tabcolsep) * \real{0.1667}}
  >{\raggedright\arraybackslash}p{(\linewidth - 10\tabcolsep) * \real{0.1667}}
  >{\raggedright\arraybackslash}p{(\linewidth - 10\tabcolsep) * \real{0.1667}}@{}}
\toprule\noalign{}
\begin{minipage}[b]{\linewidth}\raggedright
Retriever
\end{minipage} & \begin{minipage}[b]{\linewidth}\raggedright
Defense
\end{minipage} & \begin{minipage}[b]{\linewidth}\raggedright
\(\varepsilon\)
\end{minipage} & \begin{minipage}[b]{\linewidth}\raggedright
\texttt{drift\_tv} (R1)
\end{minipage} & \begin{minipage}[b]{\linewidth}\raggedright
\texttt{stance\_gap} (R2)
\end{minipage} & \begin{minipage}[b]{\linewidth}\raggedright
\texttt{poison@k}
\end{minipage} \\
\midrule\noalign{}

\bottomrule\noalign{}

BM25 & vanilla (none) & --- & 0.1500 & 0.6250 & \textbf{0.750} \\
BM25 & multi-query consistency & --- & 0.1531 & 0.6157 & \textbf{0.750} \\
BM25 & \texttt{repr\_group} (R1 only) & 1.0 & 0.1500 & 0.6250 & \textbf{0.750} \\
BM25 & \texttt{repr\_group} (R1 only) & 0.5 & 0.1500 & 0.6250 & \textbf{0.750} \\
BM25 & \texttt{repr\_group} (R1 only) & 0.25 & 0.1500 & 0.6250 & \textbf{0.750} \\
BM25 & \texttt{repr\_group} (R1 only) & 0.1 & 0.1625 & 0.6250 & \textbf{0.750} \\
BM25 & \texttt{repr\_group} (R1 only) & \textbf{0.0} & 0.1812 & 0.6250 & \textbf{0.750} \\
dense & vanilla (none) & --- & 0.1313 & 1.0000 & \textbf{1.000} \\
dense & \texttt{repr\_group} (R1 only) & 1.0 & 0.1313 & 1.0000 & \textbf{1.000} \\
dense & \texttt{repr\_group} (R1 only) & \textbf{0.0} & 0.2438 & 1.0000 & \textbf{1.000} \\
\end{longtable}
}

\textbf{Finding 2. The R1-only defense is not weak but inert.} \texttt{repr\_group} attains \emph{exactly} the same \texttt{poison@k} as \texttt{vanilla} --- 0.750 on BM25, 1.000 on dense --- at \textbf{every} budget setting, including \(\epsilon = 0\), the strictest possible constraint. The inclusion rate is not merely unchanged at the default setting; across a sweep of five values spanning the entire range of the budget parameter it does not vary at all. Tightening the budget does change R1 drift (0.1500 \(\rightarrow\) 0.1812, i.e.~it can make R1 \emph{worse}, by forcing the selection away from the natural ranking) while leaving both the R2 gap and adversarial inclusion untouched. This is the signature of a constraint acting on quantities the attack does not perturb: the budget has no purchase on the attack.

\textbf{This is a null result, so we state its strength rather than leaving a reader to guess.} A point estimate of zero can mean \texttt{}no effect'' or \texttt{}not enough data'', and the distinction matters here because the claim is that a published defense family does nothing. We therefore report paired tests on the per-query data (Table B1) instead of aggregate means:

\begin{itemize}
\tightlist
\item
  For every R1-constrained configuration, the \textbf{per-query difference vector is identically zero} --- not small, zero. The paired bootstrap CI collapses to {[}0, 0{]} and the permutation test is degenerate by construction. The informative statistic is the equivalence bound: because the differences are exactly zero rather than approximately zero, \textbf{these cells exclude any effect on adversarial-passage inclusion, not merely an effect above a threshold.}
\item
  The reason is a \textbf{ceiling}, and this is stronger than a failure to detect an effect. The unconstrained baseline already admits adversarial passages on every query it can --- 1.000 for the dense retriever, 0.750 for BM25 --- so there is no headroom for a constraint to demonstrate a reduction even in principle.
\item
  The ceiling is not an artefact of the constraint being unable to change the selection. It \emph{does} change it: on BM25 the R2 stance gap moves from 0.7436 (no defense) to 0.5385 under the R2-only constraint, so different passages are being selected. \textbf{The selection changes and not one adversarial passage is displaced.}
\end{itemize}

We give the full paired-test table, including the equivalence bounds and the contrast cells where the constraint does change the selection, in Appendix B.

This is the central empirical claim of the paper, and it applies directly to the published defense family: {[}5{]} optimises R1 composition under a fairness constraint, {[}6{]} adjusts R1 proportions and ordering, {[}7{]} controls embedder group balance, {[}8{]} equalises item-side exposure. Each constrains aggregate composition along axes the attack leaves clean, so each inherits this insensitivity. We note that multi-query consistency {[}4{]}, which operates on retrieval stability rather than composition, is also inert here (0.750).

The adaptive sweep is in Appendix C, Table C1; the ranking inverts as described below.

The adaptive attacker defeats every defense. Off-manifold filtering is by a wide margin the strongest at \(\lambda = 0\) (0.250, against 0.750 for no defense) and \textbf{collapses monotonically to 1.000 by \(\lambda = 2\)} --- a clean demonstration that its advantage is specific to the non-adaptive attacker and does not survive an adversary who observes its penalty. The representation-conserving defenses, which never achieve better than 0.625 even non-adaptively, saturate at 1.000 by \(\lambda = 0.5\): they are the most fragile under adaptation as well as the least effective without it.

\emph{Correction.} An earlier draft of this table reported 0.125 / 0.438 / 0.750 for off-manifold filtering, and described its static advantage as a factor of six over no defense. Those values came from an exploratory run written to a scratch directory excluded from version control, not from the run that produces every other number in this paper; re-running the committed configuration reproduces the values above with \textbf{18 of 18 adaptive cells bit-identical} and does not reproduce the scratch values. The qualitative conclusion is unchanged, but the static advantage is a factor of three, not six. The same error appeared in the companion paper's Table 1 and is corrected there as well; we record it here rather than silently overwriting, because a robustness factor is the kind of number a reader may quote.

\subsection{The R2 constraint works --- on the R2 metric --- and does not defend}\label{the-r2-constraint-works-on-the-r2-metric-and-does-not-defend}

Adding the cross-group R2 constraint substantially improves the R2 statistic. On BM25 the \texttt{stance\_gap} falls from 0.6250 (no defense, and R1-only defense) to \textbf{0.1429} under \texttt{repr\_both} at \(\epsilon = 1.0\) --- a \textbf{77\% reduction}, and to 0.3542 at \(\epsilon = 0.25\). The constraint is doing exactly what it was designed to do.

\textbf{Table 4.} R2 constraint effect vs.~adversarial inclusion (BM25, template + projection).

{\def\LTcaptype{none} % do not increment counter
\begin{longtable}[]{@{}
  >{\raggedright\arraybackslash}p{(\linewidth - 6\tabcolsep) * \real{0.2500}}
  >{\raggedright\arraybackslash}p{(\linewidth - 6\tabcolsep) * \real{0.2500}}
  >{\raggedright\arraybackslash}p{(\linewidth - 6\tabcolsep) * \real{0.2500}}
  >{\raggedright\arraybackslash}p{(\linewidth - 6\tabcolsep) * \real{0.2500}}@{}}
\toprule\noalign{}
\begin{minipage}[b]{\linewidth}\raggedright
Defense
\end{minipage} & \begin{minipage}[b]{\linewidth}\raggedright
\(\varepsilon\)
\end{minipage} & \begin{minipage}[b]{\linewidth}\raggedright
\texttt{stance\_gap} (R2, lower is better)
\end{minipage} & \begin{minipage}[b]{\linewidth}\raggedright
\texttt{poison@k} (security, lower is better)
\end{minipage} \\
\midrule\noalign{}

\bottomrule\noalign{}

vanilla & --- & 0.6250 & 0.750 \\
\texttt{repr\_group} (R1 only) & 0.0 & 0.6250 & 0.750 \\
\texttt{repr\_stance} (R2 only) & 1.0 & 0.4444 & 0.750 \\
\texttt{repr\_stance} (R2 only) & 0.0 & 0.5278 & 0.750 \\
\textbf{\texttt{repr\_both} (R1+R2)} & \textbf{1.0} & \textbf{0.1429} & \textbf{0.750} \\
\texttt{repr\_both} (R1+R2) & 0.25 & 0.3542 & 0.750 \\
\texttt{repr\_both} (R1+R2) & 0.0 & 0.5278 & 0.750 \\
\end{longtable}
}

\textbf{And \texttt{poison@k} is unchanged at 0.750 in every single row.} On the dense back-end the R2 constraint fails even on its own metric (\(\text{stance\_gap} \ge 0.75\) throughout, \texttt{poison@k} \(= 1.000\) at every setting).

The mechanism is visible in the implementation. The repair pass of §4.3 restores a deficient group's favourable count by swapping its weakest selected unfavourable passage for the best-scoring unselected favourable passage of that group. It improves the stance profile \emph{of the selection} --- but the substituted passage need not be legitimate. An injected passage that is favourable toward the deficient group is an admissible substitute, and because injected passages are engineered to out-rank legitimate ones, they are precisely the passages the swap selects. The defense improves its fairness statistic by substituting one adversarial passage for another.

\textbf{Finding 3.} On the controlled corpus, a distribution constraint extended to a second dimension can improve its own aggregate statistic by 77\% while adversarial inclusion does not move by a single query. On the natural corpus the same constraint \emph{does} reduce adversarial inclusion (§5.4), so Finding 3 is corpus-dependent and we do not claim it universally. What is corpus-independent is the direction of the result: \textbf{the R2 dimension is the only one along which either metric moves at all}, whereas the R1 dimension is inert in both settings. We suggest that evaluations in this area report adversarial inclusion alongside fairness metrics, because the two demonstrably move independently --- and, on the controlled corpus, in opposite directions.

\subsection{Replication on a natural corpus (BBQ)}\label{replication-on-a-natural-corpus-bbq}

The results above use a controlled corpus whose passages are template-instantiated, which is a threat to external validity. We therefore replicate on \textbf{BBQ} {[}15{]}, a naturally written bias benchmark, using BBQ's own stereotype annotations to derive stance labels (§4.1) so that no new annotation is introduced. The corpus contains 17,792 passages over 144 group-neutral queries spanning four categories (gender, disability, age, race/ethnicity).

\textbf{A note on the injection parameter, because our first attempt was wrong.} An earlier version of this experiment injected a fixed six passages per stratum and concluded that the text-only attack \texttt{}does not transfer'' to natural data. That conclusion was an artefact of the parameter, not a property of the corpus: six passages is roughly 10\% of a small candidate pool but \textasciitilde0.03\% of an 18k-passage corpus, so a fixed count measures corpus size rather than attack strength. Poisoning studies standardly report an \emph{injection rate}; we now do the same, sweeping \(\rho \in \{0.1\%, 0.25\%, 0.5\%, 1\%, 2\%\}\) of corpus size. All results below are at matched \(\rho\).

\textbf{Table 5.} Attack effect vs.~injection rate \(\rho\); no defense; \texttt{poison@k}.

\emph{Controlled corpus (1,824 passages):}

{\def\LTcaptype{none} % do not increment counter
\begin{longtable}[]{@{}
  >{\raggedright\arraybackslash}p{(\linewidth - 10\tabcolsep) * \real{0.1667}}
  >{\raggedright\arraybackslash}p{(\linewidth - 10\tabcolsep) * \real{0.1667}}
  >{\raggedright\arraybackslash}p{(\linewidth - 10\tabcolsep) * \real{0.1667}}
  >{\raggedright\arraybackslash}p{(\linewidth - 10\tabcolsep) * \real{0.1667}}
  >{\raggedright\arraybackslash}p{(\linewidth - 10\tabcolsep) * \real{0.1667}}
  >{\raggedright\arraybackslash}p{(\linewidth - 10\tabcolsep) * \real{0.1667}}@{}}
\toprule\noalign{}
\begin{minipage}[b]{\linewidth}\raggedright
\(\rho\)
\end{minipage} & \begin{minipage}[b]{\linewidth}\raggedright
passages/stratum
\end{minipage} & \begin{minipage}[b]{\linewidth}\raggedright
BM25 \texttt{poison@k}
\end{minipage} & \begin{minipage}[b]{\linewidth}\raggedright
BM25 \texttt{stance\_gap}
\end{minipage} & \begin{minipage}[b]{\linewidth}\raggedright
dense \texttt{poison@k}
\end{minipage} & \begin{minipage}[b]{\linewidth}\raggedright
dense \texttt{stance\_gap}
\end{minipage} \\
\midrule\noalign{}

\bottomrule\noalign{}

clean & 0 & --- & \textbf{0.0000} & --- & \textbf{0.0000} \\
0.1\% & 2 & \textbf{0.6875} & 0.5417 & 0.4375 & 0.2833 \\
0.25\% & 5 & 0.750 & 0.6250 & 0.625 & 0.5000 \\
1\% & 18 & 0.750 & 0.7692 & 0.625 & 0.6111 \\
2\% & 36 & 0.750 & 0.7000 & 0.625 & 0.6667 \\
\end{longtable}
}

\emph{BBQ (17,792 passages):}

{\def\LTcaptype{none} % do not increment counter
\begin{longtable}[]{@{}
  >{\raggedright\arraybackslash}p{(\linewidth - 10\tabcolsep) * \real{0.1667}}
  >{\raggedright\arraybackslash}p{(\linewidth - 10\tabcolsep) * \real{0.1667}}
  >{\raggedright\arraybackslash}p{(\linewidth - 10\tabcolsep) * \real{0.1667}}
  >{\raggedright\arraybackslash}p{(\linewidth - 10\tabcolsep) * \real{0.1667}}
  >{\raggedright\arraybackslash}p{(\linewidth - 10\tabcolsep) * \real{0.1667}}
  >{\raggedright\arraybackslash}p{(\linewidth - 10\tabcolsep) * \real{0.1667}}@{}}
\toprule\noalign{}
\begin{minipage}[b]{\linewidth}\raggedright
\(\rho\)
\end{minipage} & \begin{minipage}[b]{\linewidth}\raggedright
passages/stratum
\end{minipage} & \begin{minipage}[b]{\linewidth}\raggedright
BM25 \texttt{poison@k}
\end{minipage} & \begin{minipage}[b]{\linewidth}\raggedright
dense \texttt{poison@k}
\end{minipage} & \begin{minipage}[b]{\linewidth}\raggedright
dense \texttt{stance\_gap}
\end{minipage} & \begin{minipage}[b]{\linewidth}\raggedright
dense \texttt{in\_pool\_rate}
\end{minipage} \\
\midrule\noalign{}

\bottomrule\noalign{}

clean & 0 & --- & --- & 0.9326 & 0.5569 \\
0.1\% & 18 & \textbf{0.000} & \textbf{0.903} & 0.9395 & 0.2000 \\
0.25\% & 44 & 0.000 & 0.972 & 0.9389 & 0.1278 \\
0.5\% & 89 & 0.000 & 0.986 & 0.9696 & 0.0597 \\
1\% & 178 & 0.000 & 0.993 & 0.9942 & 0.0250 \\
2\% & 356 & 0.000 & \textbf{1.000} & 0.9958 & 0.0111 \\
\end{longtable}
}

\textbf{Finding 4. The text-only attack's effectiveness depends on how templated the corpus is, not only on injection rate.} On the controlled corpus, \textbf{8 injected passages (0.1\% of the corpus) already reach 68.75\% inclusion against BM25} --- the same order as the five-document attacks reported in the factual-poisoning literature {[}2{]}. On BBQ the same attack reaches \textbf{0.000 at every rate up to 2\%} (356 passages per stratum). The mechanism is lexical: the attack raises its score by aligning injected text with query vocabulary, and that advantage requires the injected passage to \emph{stand out} against the corpus's own vocabulary distribution. A template-instantiated corpus is repetitive, so the injected tokens dominate; naturally varied text dilutes them, and no injection rate we tested overcame the dilution. We stress that the sweep was necessary to see this --- at a fixed count we would have reported the opposite conclusion for the wrong reason.

\textbf{Finding 5. The projection attack is robust to corpus naturalness.} On the dense back-end, \texttt{poison@k} reaches \textbf{0.903 at 0.1\%} on BBQ and 1.000 by 2\%, with the utility proxy falling from 0.557 to 0.011. The projection operates in embedding space, where the perturbation's advantage does not depend on lexical distinctiveness, so it transfers where the text attack does not. Across both corpora and both injection modes, the projection attack is the only one that is consistently effective.

\textbf{Table 6.} Defense comparison on BBQ under template + projection (dense), by injection rate. \texttt{poison@k}.

{\def\LTcaptype{none} % do not increment counter
\begin{longtable}[]{@{}
  >{\raggedright\arraybackslash}p{(\linewidth - 6\tabcolsep) * \real{0.2500}}
  >{\raggedright\arraybackslash}p{(\linewidth - 6\tabcolsep) * \real{0.2500}}
  >{\raggedright\arraybackslash}p{(\linewidth - 6\tabcolsep) * \real{0.2500}}
  >{\raggedright\arraybackslash}p{(\linewidth - 6\tabcolsep) * \real{0.2500}}@{}}
\toprule\noalign{}
\begin{minipage}[b]{\linewidth}\raggedright
\(\rho\)
\end{minipage} & \begin{minipage}[b]{\linewidth}\raggedright
vanilla
\end{minipage} & \begin{minipage}[b]{\linewidth}\raggedright
\texttt{repr\_group} (R1 only) \(\varepsilon\)=0
\end{minipage} & \begin{minipage}[b]{\linewidth}\raggedright
\texttt{repr\_both} (R1+R2) \(\varepsilon\)=1.0
\end{minipage} \\
\midrule\noalign{}

\bottomrule\noalign{}

0.1\% & 0.903 & 0.903 & \textbf{0.903} \\
0.25\% & 0.972 & 0.972 & 0.951 \\
0.5\% & 0.986 & 0.986 & 0.972 \\
1\% & 0.993 & 0.993 & 0.972 \\
2\% & 1.000 & 1.000 & 0.993 \\
\end{longtable}
}

\textbf{The R2 constraint reduces adversarial inclusion on BBQ while the R1 constraint does not, and the pattern holds across the whole rate sweep.} \texttt{repr\_group} is identical to no defense at every rate (0.903 / 0.972 / 0.986 / 0.993 / 1.000), whereas \texttt{repr\_both} is weakly better at every rate, with its absolute advantage growing as the attack strengthens (0.021 at 0.5\%, 0.007 at 2\%) while its R2 advantage is much larger. The effect sizes here are smaller than the fixed-count run reported, and the reason is now clear: that run's comparison was confounded by the injection budget.

\textbf{Table 7.} R2 constraint effect on BBQ (dense, template + projection, \(\rho = 0.5\%\)).

{\def\LTcaptype{none} % do not increment counter
\begin{longtable}[]{@{}llll@{}}
\toprule\noalign{}
Defense & \(\varepsilon\) & \texttt{stance\_gap} (R2) & \texttt{poison@k} \\
\midrule\noalign{}

\bottomrule\noalign{}

vanilla & --- & 0.9696 & 0.986 \\
\texttt{repr\_group} (R1 only) & 0.0 & 0.9515 & 0.986 \\
\textbf{\texttt{repr\_both} (R1+R2)} & \textbf{1.0} & \textbf{0.5250} & \textbf{0.972} \\
\end{longtable}
}

The R2 constraint roughly halves the stance gap (0.9696 \(\rightarrow\) 0.5250) at this operating point while the R1 constraint leaves it essentially unchanged (0.9515). This is the clearest form of the paper's central claim on natural data: \textbf{the R1 dimension is inert and the R2 dimension is where the effect lives.}

We also note a ceiling effect. On BBQ the \emph{clean} retrieval already has \texttt{stance\_gap} = 0.9326, because BBQ is a bias benchmark whose contexts express stereotyped views by construction, so a query retrieves largely one-sided passages regardless of any adversary. The absolute value of \(\Delta_{\text{R2}}\) is therefore uninformative on this corpus and we rely on the change from the clean baseline throughout.

\textbf{Revised summary of findings across corpora and injection rates.}

{\def\LTcaptype{none} % do not increment counter
\begin{longtable}[]{@{}
  >{\raggedright\arraybackslash}p{(\linewidth - 6\tabcolsep) * \real{0.2500}}
  >{\raggedright\arraybackslash}p{(\linewidth - 6\tabcolsep) * \real{0.2500}}
  >{\raggedright\arraybackslash}p{(\linewidth - 6\tabcolsep) * \real{0.2500}}
  >{\raggedright\arraybackslash}p{(\linewidth - 6\tabcolsep) * \real{0.2500}}@{}}
\toprule\noalign{}
\begin{minipage}[b]{\linewidth}\raggedright
Finding
\end{minipage} & \begin{minipage}[b]{\linewidth}\raggedright
Controlled
\end{minipage} & \begin{minipage}[b]{\linewidth}\raggedright
BBQ
\end{minipage} & \begin{minipage}[b]{\linewidth}\raggedright
Status
\end{minipage} \\
\midrule\noalign{}

\bottomrule\noalign{}

R1 drift is small and non-diagnostic & yes & yes & \textbf{retained} \\
R1-only defense is inert (identical to no defense at every \(\varepsilon\)) & yes & yes & \textbf{retained} \\
R2 constraint is the only one that improves either metric & yes & yes & \textbf{retained} \\
Projection attack transfers to natural text & yes & yes (0.903 at 0.1\%) & \textbf{retained} \\
Text-only attack transfers to natural text & yes & no (0.000 at all \(\rho\)) & \textbf{corpus-dependent} (Finding 4) \\
R2 constraint reduces adversarial inclusion & no (fixed count) & yes (0.986 \(\rightarrow\) 0.972) & \textbf{rate-dependent} \\
\end{longtable}
}

The three findings that hold on both corpora and at every injection rate are the paper's claims. The two that do not --- text-attack transfer, and whether the R2 constraint reduces inclusion --- are reported with their conditions rather than dropped, because the conditions are themselves informative: both are consequences of how much lexical room the attacker has.

\textbf{Adaptive attack on BBQ.} \texttt{off-manifold} holds \texttt{poison@k} at 0.007--0.111 up to \(\lambda = 0.5\); \texttt{repr\_conserving} at 0.049--0.319; \texttt{multi\_query} at 0.083--0.535. All saturate at 1.000 by \(\lambda = 4\); see the companion paper for a full adaptive-evaluation treatment.

\subsection{Encoder robustness: a per-checkpoint property, not a family property}\label{encoder-robustness-a-per-checkpoint-property-not-a-family-property}

The results so far are measured on our own retrievers. The encoders this literature actually uses differ substantially in how much structure an attacker can exploit, so we evaluate the attack across nine retrieval back-ends: BM25 and a self-contained feature-hashing dense retriever; the learned-sparse retriever SPLADE at two sizes; and the dense semantic encoders the compared work uses --- GTE-base {[}5{]} and its larger sibling GTE-large, Contriever {[}8{]}, and E5-base-v2 {[}4, 6{]} with E5-large-v2. We include SPLADE specifically because it separates two explanations a lexical-versus-semantic comparison cannot: it is sparse like BM25 but learned like a neural encoder. E5 receives its required \texttt{query:} / \texttt{passage:} prefixes; omitting them degrades retrieval and would distort the comparison.

\textbf{Table 8.} Projection attack (\texttt{template\_plus\_projection}), \(\rho = 0.5\%\), no defense. The scale column pairs each encoder with its larger sibling from the same family and training recipe.

{\def\LTcaptype{none} % do not increment counter
\begin{longtable}[]{@{}
  >{\raggedright\arraybackslash}p{(\linewidth - 10\tabcolsep) * \real{0.1667}}
  >{\raggedright\arraybackslash}p{(\linewidth - 10\tabcolsep) * \real{0.1667}}
  >{\raggedright\arraybackslash}p{(\linewidth - 10\tabcolsep) * \real{0.1667}}
  >{\raggedright\arraybackslash}p{(\linewidth - 10\tabcolsep) * \real{0.1667}}
  >{\raggedright\arraybackslash}p{(\linewidth - 10\tabcolsep) * \real{0.1667}}
  >{\raggedright\arraybackslash}p{(\linewidth - 10\tabcolsep) * \real{0.1667}}@{}}
\toprule\noalign{}
\begin{minipage}[b]{\linewidth}\raggedright
Retriever
\end{minipage} & \begin{minipage}[b]{\linewidth}\raggedright
representation
\end{minipage} & \begin{minipage}[b]{\linewidth}\raggedright
size
\end{minipage} & \begin{minipage}[b]{\linewidth}\raggedright
R1 drift
\end{minipage} & \begin{minipage}[b]{\linewidth}\raggedright
R2 gap
\end{minipage} & \begin{minipage}[b]{\linewidth}\raggedright
\texttt{poison@k}
\end{minipage} \\
\midrule\noalign{}

\bottomrule\noalign{}

BM25 & lexical, unsupervised sparse & --- & 0.1625 & 0.7436 & 0.750 \\
dense (feature hashing) & hashed, not semantic & --- & 0.1625 & 1.0000 & 1.000 \\
SPLADE & learned sparse & 66M & 0.2687 & 1.0000 & 0.750 \\
SPLADE large & learned sparse & 110M & 0.2687 & 1.0000 & 1.000 \\
GTE-base & dense semantic & 110M & 0.4688 & 1.0000 & 1.000 \\
\textbf{GTE-large} & dense semantic & 335M & 0.4188 & 1.0000 & 1.000 \\
Contriever & dense semantic & 110M & 0.3687 & 1.0000 & 1.000 \\
E5-base-v2 & dense semantic & 110M & 0.4000 & 1.0000 & 1.000 \\
\textbf{E5-large-v2} & dense semantic & 335M & 0.3750 & 1.0000 & 1.000 \\
\end{longtable}
}

Every back-end is attacked, and the R2 gap saturates at its maximum for all but BM25. The dense encoders' R1 drift (0.37--0.47) is \textbf{2.3\(\times\)--2.9\(\times\)} that of BM25 (0.16): the projection attack perturbs group composition \emph{more} in a continuous semantic space, the opposite of what we expected when adding this experiment.

\textbf{Table 9.} Text-only (lexical) attack, \(\rho = 0.5\%\) and \(2\%\), no defense. This is where the back-ends separate.

{\def\LTcaptype{none} % do not increment counter
\begin{longtable}[]{@{}
  >{\raggedright\arraybackslash}p{(\linewidth - 10\tabcolsep) * \real{0.1667}}
  >{\raggedright\arraybackslash}p{(\linewidth - 10\tabcolsep) * \real{0.1667}}
  >{\raggedright\arraybackslash}p{(\linewidth - 10\tabcolsep) * \real{0.1667}}
  >{\raggedright\arraybackslash}p{(\linewidth - 10\tabcolsep) * \real{0.1667}}
  >{\raggedright\arraybackslash}p{(\linewidth - 10\tabcolsep) * \real{0.1667}}
  >{\raggedright\arraybackslash}p{(\linewidth - 10\tabcolsep) * \real{0.1667}}@{}}
\toprule\noalign{}
\begin{minipage}[b]{\linewidth}\raggedright
Retriever
\end{minipage} & \begin{minipage}[b]{\linewidth}\raggedright
representation
\end{minipage} & \begin{minipage}[b]{\linewidth}\raggedright
\texttt{poison@k} @0.5\%
\end{minipage} & \begin{minipage}[b]{\linewidth}\raggedright
@2\%
\end{minipage} & \begin{minipage}[b]{\linewidth}\raggedright
R1 drift
\end{minipage} & \begin{minipage}[b]{\linewidth}\raggedright
\texttt{fav\_g2}
\end{minipage} \\
\midrule\noalign{}

\bottomrule\noalign{}

BM25 & lexical, unsupervised sparse & 0.7500 & 0.7500 & 0.1625 & 0.7436 \\
dense (feature hashing) & hashed (keeps lexical surface) & 0.6250 & 0.6250 & 0.0812 & 0.6167 \\
\textbf{SPLADE} & \textbf{learned sparse} & \textbf{0.6250} & 0.6250 & 0.2188 & \textbf{1.0000} \\
\textbf{SPLADE large} & learned sparse & \textbf{0.5625} & 0.5625 & 0.1875 & 0.8810 \\
Contriever & dense semantic & 0.5625 & 0.5625 & 0.1750 & 0.9444 \\
GTE-base & dense semantic & \textbf{0.0625} & 0.0625 & 0.0000 & n/a \\
\textbf{GTE-large} & dense semantic & \textbf{0.5000} & 0.5000 & 0.1313 & 0.9444 \\
E5-base-v2 & dense semantic & \textbf{0.0625} & 0.0625 & 0.0000 & n/a \\
\textbf{E5-large-v2} & dense semantic & \textbf{0.0000} & 0.0000 & 0.0000 & n/a \\
\end{longtable}
}

Three observations, and they are the finding.

\textbf{Finding 7. Sparsity, not \texttt{}neuralness'', predicts susceptibility at a fixed size, but it is a correlate rather than a law.} SPLADE learns its term weights --- it is a neural encoder trained on relevance --- yet it is as susceptible as BM25 (0.6250 versus 0.7500), because its representation is still a sparse bag of term weights. The attack raises a passage's score by appending query-aligned terms, which requires the representation to expose \textbf{per-term contributions}: a sparse retriever does that whether or not its weights are learned, and a dense encoder absorbs the appended text into a single vector. SPLADE is in fact the \emph{most} efficiently attacked of the nine on the R2 dimension, reaching the maximum stance gap 1.0000 where BM25 reaches 0.7436 --- its learned weighting evidently concentrates score on exactly the terms the attacker appends. But sparsity is a correlate, not a law, and two cases bound it: Contriever is dense and semantic yet nearly as susceptible as BM25 (0.5625), and GTE-large is dense and semantic yet eight times more susceptible than its own base checkpoint.

\textbf{Finding 11. Susceptibility is not monotone in encoder scale, and scaling can destroy resistance entirely.} Table 9's clearest result is unexpected: GTE-base and E5-base-v2 are equally resistant at base size (both 0.0625 --- the attack succeeds on three of 48 queries), and enlarging both by the same factor of three moves them in \emph{opposite} directions. \textbf{E5-large-v2 becomes more resistant (0.0625 \(\rightarrow\) 0.0000, the attack now fails on every query); GTE-large becomes eight times more susceptible (0.0625 \(\rightarrow\) 0.5000, succeeding on 24 of 48).} SPLADE is susceptible at both sizes with little change (0.6250 \(\rightarrow\) 0.5625). Per-query inclusion is binary, so this is a shift in \emph{how many} queries are compromised, not a drift in a continuous score.

\textbf{The mechanism is not geometry.} We tested the obvious explanation and it is wrong, which makes the result more informative. If GTE-large's space were more anisotropic --- vectors collapsed into a narrower cone --- a fixed textual nudge would move rank further, and the finding would reduce to a known property. Measured, GTE-base and GTE-large are near-identical in mean pairwise cosine (0.8484 vs 0.8579), effective dimensionality (13.2 both, against 768 and 1024 nominal), and spread of query similarity (0.0860 vs 0.0895); the E5 pair behaves the same (0.8291 / 0.8375; 13.9 / 14.8). What differs is the gain the attack buys:

\textbf{Table 10.} Where the GTE scale effect comes from. 36 injected passages, 48 queries, no defense, \(\rho = 0.5\%\).

{\def\LTcaptype{none} % do not increment counter
\begin{longtable}[]{@{}
  >{\raggedright\arraybackslash}p{(\linewidth - 4\tabcolsep) * \real{0.3333}}
  >{\raggedright\arraybackslash}p{(\linewidth - 4\tabcolsep) * \real{0.3333}}
  >{\raggedright\arraybackslash}p{(\linewidth - 4\tabcolsep) * \real{0.3333}}@{}}
\toprule\noalign{}
\begin{minipage}[b]{\linewidth}\raggedright
Quantity
\end{minipage} & \begin{minipage}[b]{\linewidth}\raggedright
GTE-base
\end{minipage} & \begin{minipage}[b]{\linewidth}\raggedright
GTE-large
\end{minipage} \\
\midrule\noalign{}

\bottomrule\noalign{}

mean cosine gain of poison over clean & \textbf{+0.0115} & \textbf{+0.0200} \\
mean best-poison cosine & 0.8907 & 0.8975 \\
mean top-5 clean threshold & 0.8982 & 0.8988 \\
mean margin (best poison \(-\) threshold) & \textbf{\$-\(0.0075** | **\)-\$0.0013} & \\
\texttt{poison@k} & 0.0625 & 0.5000 \\
\end{longtable}
}

The two encoders place the same injected text at almost the same distance from the query (0.8907 vs 0.8975), and the legitimate competition sits at the same threshold (0.8982 vs 0.8988). The difference is that \textbf{the identical appended vocabulary buys 74\% more similarity on GTE-large (+0.0200 vs +0.0115)}. Because the mean margin is only about \$-\$0.007, a gain difference of +0.0085 carries a large fraction of queries across the threshold --- exactly the 3-to-24 shift. Susceptibility is the interaction between the attack's vocabulary and the encoder's learned weighting of it, and none of the coarse descriptors we tried --- architecture, dimensionality, size, geometry, sparsity --- predicts it.

\textbf{Table 11.} R1-only constraint inertness across encoders and scales. \texttt{template\_plus\_projection}, \(\rho = 0.5\%\); \(\Delta\) = adversarial-passage inclusion minus no defense.

{\def\LTcaptype{none} % do not increment counter
\begin{longtable}[]{@{}
  >{\raggedright\arraybackslash}p{(\linewidth - 6\tabcolsep) * \real{0.2500}}
  >{\raggedright\arraybackslash}p{(\linewidth - 6\tabcolsep) * \real{0.2500}}
  >{\raggedright\arraybackslash}p{(\linewidth - 6\tabcolsep) * \real{0.2500}}
  >{\raggedright\arraybackslash}p{(\linewidth - 6\tabcolsep) * \real{0.2500}}@{}}
\toprule\noalign{}
\begin{minipage}[b]{\linewidth}\raggedright
Retriever
\end{minipage} & \begin{minipage}[b]{\linewidth}\raggedright
\texttt{repr\_group} \(\varepsilon\)=0
\end{minipage} & \begin{minipage}[b]{\linewidth}\raggedright
\texttt{repr\_both} \(\varepsilon\)=0
\end{minipage} & \begin{minipage}[b]{\linewidth}\raggedright
\(\Delta\)
\end{minipage} \\
\midrule\noalign{}

\bottomrule\noalign{}

BM25 & 0.7500 & 0.7500 & \textbf{0.0000} \\
dense (hash) & 1.0000 & 1.0000 & \textbf{0.0000} \\
SPLADE distil & 0.7500 & 0.7500 & \textbf{0.0000} \\
SPLADE large & 1.0000 & 1.0000 & \textbf{0.0000} \\
GTE-base & 1.0000 & 1.0000 & \textbf{0.0000} \\
GTE-large & 1.0000 & 1.0000 & \textbf{0.0000} \\
E5-base-v2 & 1.0000 & 1.0000 & \textbf{0.0000} \\
E5-large-v2 & 1.0000 & 1.0000 & \textbf{0.0000} \\
\end{longtable}
}

\textbf{Finding 6b. The R1 constraint is inert across every encoder and at every scale tested.} In all sixteen constrained configurations the difference is exactly zero. The projection attack reaches complete inclusion on seven of the eight, so the constraint is inert in the worst case rather than a marginal one. This is the third independent setting in which the inertness holds, after the two corpora and the injection-rate sweep.

\subsection{Adaptive attacker: no defense survives an informed adversary}\label{adaptive-attacker-no-defense-survives-an-informed-adversary}

The inertness results above are measured against a \emph{static} attacker. The security literature's standard objection --- and the reason adaptive evaluation is mandatory --- is that a defense may hold against an attacker who does not know it and fail against one who does. It matters here for two independent reasons, both predicted by §6.5: our defenses are either randomised with public parameters or penalise an observable quantity.

For each defense the attacker is given the defense's own penalty function and searches over perturbation strength \(\lambda\), using the same subspace-projection construction. We also report \texttt{usage}, the mean cosine similarity between each injected passage's perturbed and original embedding, so that a defense cannot be credited for the attacker's self-inflicted degradation: if the injected material were destroyed, low adversarial inclusion would be meaningless.

\emph{The per-\(\lambda\) grid is in Appendix C, Table C1.}

Off-manifold filtering is by a wide margin the strongest statically --- it reduces adversarial inclusion from 0.750 for no defense to 0.250, a factor of three --- and \textbf{collapses monotonically to complete failure by \(\lambda = 2\)}. Its static advantage is exactly what §6.5's Corollary 2 predicts: the defense penalises vectors that leave the corpus manifold, which is precisely where the projection perturbation places them, so once the attacker trades perturbation strength against penalty, the penalty becomes the attacker's constraint. The two composition-constrained defenses, which never beat 0.625 even statically, saturate by \(\lambda = 0.5\). \textbf{The static ranking of defenses inverts under adaptation, and a defense selected on its static advantage is defeated as comprehensively as one selected on structural grounds.}

\textbf{The attacker pays almost nothing for this.} \texttt{usage} stays at 0.945--0.979 through \(\lambda = 1\) and 0.883 at \(\lambda = 2\). At the point where the strongest defense has already lost most of its advantage, the injected material is over 94\% semantically intact. The collapse therefore cannot be attributed to the attack degrading its own passages.

\textbf{Correction.} An earlier draft reported 0.125 / 0.438 / 0.750 for off-manifold filtering and described its static advantage as a factor of six over no defense. Those values came from an exploratory run written to a scratch directory excluded from version control, not from the run that produces every other number in this paper; re-running the committed configuration reproduces Table C1 with 18 of 18 adaptive cells bit-identical. The qualitative conclusion is unchanged, but the static advantage is a factor of three, not six. We record it rather than silently overwriting, because a robustness factor is the kind of number a reader may quote.

\subsection{Generation-stage propagation: does the retrieval skew reach the output?}\label{generation-stage-propagation-does-the-retrieval-skew-reach-the-output}

Every result above is measured at the retrieval layer. That is deliberate --- it makes the mechanism measurable without a generator and reproducible on commodity hardware --- but it leaves the question a reviewer will ask: \textbf{does a retrieval-layer stance skew change what the system says?} We answer it with a fixed generator, varying only the retrieved context, so that any difference between conditions is attributable to retrieval.

\textbf{Protocol.} For each query we generate under four conditions, holding model, prompt, decoding (temperature 0) and question text fixed: \texttt{clean} (unpoisoned retrieval), \texttt{poisoned} (pairwise injection, no defence), \texttt{r1only} (R1 constraint), \texttt{r2both} (R1+R2 constraint). The primary probe is \textbf{forced choice}: the generator must select between a favourable and an unfavourable statement about a group --- the generation-layer analogue of the R2 dimension, and the same form of measurement used for attack success in the poisoning literature. Option order is randomised per (query, condition, group) so that a position bias cannot masquerade as a stance effect. We additionally collect free-form answers and a per-passage attribution probe.

Generator: DeepSeek (\texttt{deepseek-chat} at temperature 0), selected because it is reachable without a proxy from the network this work used. 207 API calls, 49.2k prompt tokens, 125 s of model time.

Read Table C2 and Table 11 together: the same experiment, on the same corpus at the same injection rate, yields a 46\% rise in the absolute gap under one instrument and a rise of 0.0008--0.0018 under the other. We keep both because that discrepancy is itself the finding.

\textbf{Finding 8. On the controlled corpus the retrieval-layer skew propagates to the generated output and replicates across three model families, but it does not survive a naturally written corpus.} We re-ran this evaluation with a continuous, entailment-based stance metric (a local NLI model scoring \texttt{P(entail\ \textbar{}\ answer,\ favourable)\ \$-\$\ P(entail\ \textbar{}\ answer,\ unfavourable)}, calibrated in §5.7.1) across four generators spanning three model families, one of them (Mistral-7B) open-weight and served locally rather than by an API, so that the result does not rest on hosted models alone. The per-group shifts are large, consistent in sign, and highly significant everywhere:

\textbf{Table 12.} Generation-layer stance, entailment-scored, four generators spanning three model families. Controlled corpus, GTE-base retrieval, \(\rho = 2\%\), 48 queries, paired permutation test against \texttt{clean}.

{\def\LTcaptype{none} % do not increment counter
\begin{longtable}[]{@{}
  >{\raggedright\arraybackslash}p{(\linewidth - 12\tabcolsep) * \real{0.1429}}
  >{\raggedright\arraybackslash}p{(\linewidth - 12\tabcolsep) * \real{0.1429}}
  >{\raggedright\arraybackslash}p{(\linewidth - 12\tabcolsep) * \real{0.1429}}
  >{\raggedright\arraybackslash}p{(\linewidth - 12\tabcolsep) * \real{0.1429}}
  >{\raggedright\arraybackslash}p{(\linewidth - 12\tabcolsep) * \real{0.1429}}
  >{\raggedright\arraybackslash}p{(\linewidth - 12\tabcolsep) * \real{0.1429}}
  >{\raggedright\arraybackslash}p{(\linewidth - 12\tabcolsep) * \real{0.1429}}@{}}
\toprule\noalign{}
\begin{minipage}[b]{\linewidth}\raggedright
Generator
\end{minipage} & \begin{minipage}[b]{\linewidth}\raggedright
\(\Delta\) gap (absolute)
\end{minipage} & \begin{minipage}[b]{\linewidth}\raggedright
p
\end{minipage} & \begin{minipage}[b]{\linewidth}\raggedright
\textbf{\(\Delta\) group 1}
\end{minipage} & \begin{minipage}[b]{\linewidth}\raggedright
p
\end{minipage} & \begin{minipage}[b]{\linewidth}\raggedright
\textbf{\(\Delta\) group 2}
\end{minipage} & \begin{minipage}[b]{\linewidth}\raggedright
p
\end{minipage} \\
\midrule\noalign{}

\bottomrule\noalign{}

qwen-turbo & +0.0008 & 0.32 & \textbf{\$-\(0.1297** | **<0.0001** | **+0.1305** | **<0.0001** |
| qwen-plus | +0.0018 | 0.042 | **\)-\$0.1280} & \textbf{0.0035} & \textbf{+0.1298} & \textbf{0.0015} \\
qwen-max & \$-\(0.0491 | 0.18 | **\)-\(0.1785** | **0.0001** | **+0.1294** | **0.0029** |
| **mistral-7b** (local, open weights) | +0.0025 | 0.42 | **\)-\$0.1297** & \textbf{0.0001} & \textbf{+0.1322} & \textbf{\textless0.0001} & & \\
\end{longtable}
}

\textbf{The absolute cross-group gap --- the metric we originally used --- is blind to this effect, and we had reproduced one level up the exact error we criticise in §5.3.} Under injection group 1's stance collapses (\$-\$0.13 to \$-\$0.18) while group 2's rises (+0.13), in every generator, at p \(\leq\) 0.0035. Because the corpus carries a pre-existing stance asymmetry, \texttt{\textbar{}stance(g1)\ \$-\$\ stance(g2)\textbar{}} moves by less than 0.002 for two of the three generators and by \$-\$0.05 (the wrong sign) for the third. Taking an absolute difference between two groups destroys precisely the signal, for the same reason it does at the retrieval layer: the attack does not make the system uniformly more skewed, it suppresses one group and elevates the other, and the gap between them is dominated by a corpus-level constant that the attack does not touch.

This vindicates the qualitative claim in Finding 9 --- the skew saturates one group rather than shifting everything --- while changing the evidence for it. Findings 8 and 9 as originally stated rested on a forced-choice probe whose original numbers (gap 0.2708 \(\rightarrow\) 0.3958, a 46\% rise) do \textbf{not} reproduce under entailment scoring: the gap delta is +0.0008 / +0.0018 / \$-\$0.0491 across the three generators. The forced-choice probe is not wrong so much as compressed: it forces a binary selection between two maximally opposed statements, which pushes both groups' rates toward ceiling (the original table showed group 2 at 0.979 \(\rightarrow\) 1.000), whereas a continuous entailment score can register the collapse of the suppressed group at 0.132 \(\rightarrow\) 0.002, which is the larger and more consistent effect. \textbf{We therefore report the per-group entailment shifts as the primary generation-layer evidence, and we demote the absolute gap to a control that fails.}

\textbf{Finding 9. The skew does not appear as a uniform shift; it suppresses one group and elevates the other.} Under injection the suppressed group's stance falls by 0.13--0.18 and the favoured group's rises by 0.13, in all three generators. The magnitude of the two shifts is nearly equal, which is why the absolute difference between them is nearly unchanged: the attack relocates both groups rather than opening a gap between them. This is visible only because we measure the two groups separately, and it is the generation-layer counterpart of the R1/R2 distinction the paper is built on.

\textbf{Finding 10. Neither constraint improves the generation layer.} \texttt{r1only} leaves both per-group shifts identical to no defence to four decimal places (\(\Delta\)g1 = \$-\$0.1297, \(\Delta\)g2 = +0.1305 for qwen-turbo, identical values), and \texttt{r2both} moves them by less than 0.002. This is consistent with the retrieval-layer result: the input to the generator is unchanged, so the output is unchanged. \textbf{The constraints we proposed, evaluated here, do not repair the downstream harm.}

\textbf{5.7.1 Why the natural corpus appeared to show nothing: the probe, not the corpus.} The free-form measurement above is valid only where answers actually take a position, and on a naturally written corpus they largely do not. Measuring the \emph{commitment} of each answer --- the larger of its two entailment probabilities, \texttt{max(P(fav),\ P(unfav))} --- separates a real null from an instrument that cannot register anything:

\textbf{Table 13.} Answer commitment by corpus, clean condition. An answer with commitment near zero takes no position, so no stance change can appear in it however large the retrieval-layer skew.

{\def\LTcaptype{none} % do not increment counter
\begin{longtable}[]{@{}lll@{}}
\toprule\noalign{}
Corpus & mean commitment & answers below 0.1 \\
\midrule\noalign{}

\bottomrule\noalign{}

controlled & 0.1555 & 84.4\% \\
\textbf{BBQ (full)} & \textbf{0.0071} & \textbf{99.0\%} \\
\end{longtable}
}

The gap is 22\(\times\), and it means the free-form statistic on BBQ is near zero under \emph{every} condition: the effect we were looking for could not have appeared there. This also bounds the controlled-corpus result more honestly than we previously stated it --- 84.4\% of those answers are non-committal too, so the shift reported above rests on roughly a sixth of them.

\textbf{A probe that forces commitment detects the effect on BBQ.} We therefore re-ran the generation stage with the forced-choice probe used in earlier drafts of this work: the generator must choose between a favourable and an unfavourable statement about the group, drawn from the same template inventory the attack manipulates, with option order randomised per item so a position bias cannot masquerade as a stance effect. Before using it we verified it carries signal rather than position: the reply was a parseable choice in 100\% of cases, and the probability of choosing position A is 100\% when the favourable statement sits at A against 35\% when it sits at B --- a 65-point swing, so the choice tracks content.

\textbf{Table 14.} Forced-choice stance on the full BBQ corpus, 144 queries, per-group favourable-choice rate, paired permutation test against \texttt{clean}. Only shifts significant at \(p < 0.05\) are listed.

{\def\LTcaptype{none} % do not increment counter
\begin{longtable}[]{@{}llll@{}}
\toprule\noalign{}
Generator & Group & \(\Delta\) favourable rate & p \\
\midrule\noalign{}

\bottomrule\noalign{}

\textbf{qwen-turbo} & woman & \textbf{+0.6333} & \textbf{\textless0.0001} \\
\textbf{qwen-turbo} & nondisabled & \textbf{+0.5714} & \textbf{\textless0.0001} \\
\textbf{qwen-turbo} & man & \textbf{+0.5000} & \textbf{0.0003} \\
\textbf{qwen-turbo} & young & +0.2449 & 0.0121 \\
\textbf{qwen-turbo} & disabled & +0.2381 & 0.0274 \\
qwen-plus & man & +0.2333 & 0.0379 \\
mistral-7b & --- & none reached significance & --- \\
\end{longtable}
}

\textbf{Finding 8, restated.} The retrieval-layer skew does reach the generated output on a naturally written corpus. What failed on BBQ was the free-form probe, which could not have detected the effect because 99.0\% of the answers it scores take no position at all. With a probe that forces commitment, one generator shows five per-group shifts at \(p < 0.05\) with the largest at \(+0.63\).

Three qualifications belong with that, and we state them rather than leaving them to be found:

\begin{itemize}
\tightlist
\item
  \textbf{The effect is group-specific, not uniform.} On BBQ it concentrates in woman, nondisabled, man, young and disabled; on the controlled corpus (Table 12) it concentrates in asian, elderly and white. Averaging over groups, which is what a per-group \emph{mean} does and what our earlier free-form analysis did, understates it substantially.
\item
  \textbf{It is generator-dependent.} qwen-turbo carries most of the signal on both corpora; mistral-7b shows none that reaches significance on BBQ, and the one API generator we ran at largest scale shows a single group. A single dense number such as \texttt{}the gap rose 46\%'' is therefore the wrong form for this result, and we no longer report one.
\item
  \textbf{Concentration is not the same as a large average effect.} Five significant groups out of eight, each with a large rate shift, is a stronger statement than the average gap suggests, and a weaker one than \texttt{}the system's stance changed across the board''.
\end{itemize}

\textbf{This is the third instrument failure in the paper, and it belongs with the other two.} The free-form probe joins the reference-based stance statistics and the absolute cross-group gap as a measurement that returns a plausible number while being structurally unable to see the thing being measured. Its signature is the one we have now seen three times: a statistic that comes out near zero for reasons that have nothing to do with the phenomenon. We report it in the results rather than the limitations because it changes a conclusion we had already drawn.

\textbf{5.7.2 Two metric failures, and what they were caused by.}

\textbf{5.7.2 Two metric failures, and what they were caused by.} Both are worth recording because each cost us a wrong result before being found.

\emph{Stance direction.} The first implementation scored stance as \texttt{P(entail\ \textbar{}\ statement,\ answer)}. Calibrated against the correct convention on the same data, the reversed form returns ±0.002 for a favourable answer, an unfavourable answer, an unrelated answer and a vacuous one alike: it cannot separate anything, and it fails silently by returning zeros rather than an obvious error. The correct direction, \texttt{P(entail\ \textbar{}\ answer,\ statement)}, returns +0.997 and \$-\$0.997 for the two extremes. The convention is not interchangeable and the failure mode is indistinguishable from a genuine null result --- a mistake worth flagging for anyone building this measurement.

\emph{Attribution direction.} The expected-attributed-exposure (EAE-D) statistic of Kim \& Diaz {[}8{]} was implemented by asking the generator, per retrieved passage, whether its answer relied on that passage. It returned \textbf{1.000 in every condition}, which we previously reported as a failed metric. The cause is again direction: attribution must ask whether the \textbf{passage entails the answer}, not the reverse. A passage the answer used scores 0.996 in the correct direction and 0.008 in the reversed one; an unused passage scores 0.000 in both. With the direction fixed the statistic discriminates properly --- 0.798--0.828 clean, rising to 0.831--0.968 under injection, since poisoned contexts are engineered to be the passages the answer leans on. The earlier self-report probe measured the judge's agreeableness, not the answer's provenance.

\textbf{Scope.} Three generators from two model families, one retrieval backbone (GTE-base), one injection rate, and 48 retrieval slots that correspond to only 16 distinct question texts (§8, Limitation 6). The per-group shifts are consistent in sign and significance across all three generators, which is stronger evidence than the single-generator version this section previously reported; the magnitudes, however, should not be extrapolated to generators we did not run.

\begin{center}\rule{0.5\linewidth}{0.5pt}\end{center}

\section{Why Distribution-Level Defense Has a Limit}\label{why-distribution-level-defense-has-a-limit}

Findings 2, 3 and 6 are not tuning failures. They follow from what a distributional constraint \emph{is}, and we give the derivation here.

\subsection{Setup}\label{setup}

Let \(\mathcal{C}\) be the corpus; each passage \(d\) carries a group \(g(d) \in \mathcal{G} \cup \{\bot\}\) and a stance \(s(d) \in \{+1,-1,0\}\) (\(+1\) favourable toward \(g(d)\)). Let \(\mathcal{A} \subseteq \mathcal{C}\) be the passages the attacker controls, \(n = |\mathcal{A}|\). A defense selects \(\mathcal{D}_q \subseteq \mathcal{C}\) with \(|\mathcal{D}_q| = k\) (we write \(\mathcal{D}\) where \(q\) is clear). Define

\[p_{\text{grp}}(g;\mathcal{D}) = \frac{|\{d \in \mathcal{D} : g(d) = g,\ s(d)\neq 0\}|}{|\{d \in \mathcal{D} : g(d)\neq\bot,\ s(d)\neq 0\}|}, \qquad p_{\text{fav}}(g;\mathcal{D}) = \frac{|\{d \in \mathcal{D} : g(d) = g,\ s(d)=+1\}|}{|\{d \in \mathcal{D} : g(d) = g,\ s(d)\neq 0\}|}.\]

\textbf{Definition (\(\varepsilon\)-feasible).} Fix references \(p^{*}_{\text{grp}}, p^{*}_{\text{fav}}\). A selection \(\mathcal{D}\) is \emph{\(\varepsilon\)-feasible}, written \(\mathcal{D} \in \mathcal{R}_\varepsilon\), if for every \(g \in \mathcal{G}\)

\[\bigl|p_{\text{grp}}(g;\mathcal{D}) - p^{*}_{\text{grp}}(g)\bigr| \le \varepsilon \quad\text{and}\quad \bigl|p_{\text{fav}}(g;\mathcal{D}) - p^{*}_{\text{fav}}(g)\bigr| \le \varepsilon. \tag{1}\]

A distribution-constrained defense accepts exactly \(\mathcal{R}_\varepsilon\). All defenses in the family we evaluate --- the R1 constraint of §4.3, and by inspection the objective of {[}5{]} and the proportion adjustment of {[}6{]} --- are of this form, with \(\mathcal{R}\) built from \emph{clean} statistics.

\subsection{Proposition 1 --- the injection is admissible at every tolerance}\label{proposition-1-the-injection-is-admissible-at-every-tolerance}

\textbf{Proposition 1.} Let the attacker inject \(n\) passages that are (i) group-relevant, (ii) split between groups in proportion \(p^{*}_{\text{grp}}\), and (iii) stance-balanced so that \(p_{\text{fav}}(g) = p^{*}_{\text{fav}}(g)\) for every \(g\). Then any selection consisting only of injected passages is \(\varepsilon\)-feasible for every \(\varepsilon \ge 0\).

\emph{Proof.} Take \(\mathcal{D} \subseteq \mathcal{A}\) with \(|\mathcal{D}| = k\). By (ii), the count of passages of group \(g\) in \(\mathcal{D}\) is \(p^{*}_{\text{grp}}(g)\,k\) up to rounding, so \(p_{\text{grp}}(g;\mathcal{D}) = p^{*}_{\text{grp}}(g)\), and the first inequality of (1) holds with equality. By (iii), \(p_{\text{fav}}(g;\mathcal{D}) = p^{*}_{\text{fav}}(g)\), so the second holds with equality. Hence \(\mathcal{D} \in \mathcal{R}_\varepsilon\) for all \(\varepsilon \ge 0\). \(\square\)

\textbf{Consequence.} A defense whose acceptance test is membership in \(\mathcal{R}_\varepsilon\) admits a \textbf{fully adversarial} selection at any tolerance, including \(\varepsilon = 0\). Tightening the budget does not exclude the injection; it constrains the attacker's \emph{composition}, which the attacker chooses freely. This is the mechanism behind Findings 2 and 6, and it explains why the empirical \texttt{poison@k} is \emph{identical} to no defense at every \(\varepsilon\) rather than merely close to it: the injected selection is not near the boundary of \(\mathcal{R}_\varepsilon\), it is at the reference point.

\textbf{Empirical check.} Table C3 reports the prediction of Prop. 1 against every configuration we ran: if the injected selection sits at the reference point, then constraining the R1 budget must leave adversarial inclusion \emph{exactly} unchanged at every tolerance.

\emph{(entries are \(\Delta\)\texttt{poison@k}; 0.000 in all 24 measured cells)}

Prop. 1 predicts \(\Delta\) = 0 exactly, and \(\Delta\) = 0 exactly in every one of the 24 cells measured --- not approximately, which a tuning-failure explanation would predict. The equivalence is reproduced by \texttt{analysis/verify\_theory.py}. We state the count explicitly because the exactness is the point: the constrained selection is not \emph{near} the reference composition, it \emph{is} the reference composition.

\subsection{\texorpdfstring{6.3 Proposition 2 --- the minimum injection for an \(\varepsilon\)-shift}{6.3 Proposition 2 --- the minimum injection for an \textbackslash varepsilon-shift}}\label{proposition-2-the-minimum-injection-for-an-varepsilon-shift}

Proposition 1 shows the attacker \emph{can} be admissible; we now give the budget at which it becomes \emph{effective}.

\textbf{Proposition 2.} Fix a query and a group \(g\). Suppose the clean selection contains \(m_g \le k\) passages about \(g\), of which \(f_g\) are favourable, so \(p^{*}_{\text{fav}}(g) = f_g/m_g\). An attacker who replaces \(j \le \min(n, m_g)\) of these with injected passages about \(g\) (favourable, to flood; unfavourable, to suppress) violates the stance constraint of (1) for \(g\) if and only if

\[j > \varepsilon\, m_g, \qquad\text{so}\qquad j^{*}_g(\varepsilon) = \lceil \varepsilon\, m_g \rceil + 1. \tag{2}\]

\emph{Proof.} Flooding: \(p_{\text{fav}}(g) = (f_g + j)/m_g\) for \(j \le m_g - f_g\), so \(p_{\text{fav}}(g) - p^{*}_{\text{fav}}(g) = j/m_g > \varepsilon \iff j > \varepsilon m_g\). Suppressing: \(p_{\text{fav}}(g) = (f_g - j)/m_g\), and \(p^{*}_{\text{fav}}(g) - p_{\text{fav}}(g) = j/m_g\), giving the same condition. The smallest integer satisfying the strict inequality is \(\lceil \varepsilon m_g\rceil + 1\), which is needed to handle the case where \(\varepsilon m_g\) is an integer. \(\square\)

\textbf{Corollary 1 (the \emph{constraint} cost is at most \(k\)).} Since \(m_g \le k\), \(j^{*}_g(\varepsilon) \le k\) for all \(\varepsilon \le 1\). With \(k\) typically 3--5, violating the constraint is cheap in the currency of the constraint itself.

\textbf{Corollary 2 (suppression is cheaper than flooding when the clean retrieval is one-sided).} If \(f_g = 0\) then \(p^{*}_{\text{fav}}(g) = 0\) and the flooding direction requires no violation at all, while the suppression direction must overcome the full rate; if \(f_g = m_g\) the reverse holds. The cheaper direction is always the one pushing an already-extreme rate further toward its extreme --- consistent with the empirical Finding 9, where the attack drove one group's favourable rate to ceiling rather than shifting both.

\textbf{A correction we had to make.} An earlier version of this section concluded from Corollary 1 that \texttt{}the number of passages required is bounded by a small constant independent of corpus size''. That conflates two requirements that behave very differently:

{\def\LTcaptype{none} % do not increment counter
\begin{longtable}[]{@{}
  >{\raggedright\arraybackslash}p{(\linewidth - 4\tabcolsep) * \real{0.3333}}
  >{\raggedright\arraybackslash}p{(\linewidth - 4\tabcolsep) * \real{0.3333}}
  >{\raggedright\arraybackslash}p{(\linewidth - 4\tabcolsep) * \real{0.3333}}@{}}
\toprule\noalign{}
\begin{minipage}[b]{\linewidth}\raggedright
Requirement
\end{minipage} & \begin{minipage}[b]{\linewidth}\raggedright
Governed by
\end{minipage} & \begin{minipage}[b]{\linewidth}\raggedright
Depends on corpus size?
\end{minipage} \\
\midrule\noalign{}

\bottomrule\noalign{}

\textbf{Violate the constraint} (Prop. 2) & \(\varepsilon\) and \(m_g \le k\) & \textbf{No} \\
\textbf{Be retrieved at all} & whether injected passages out-rank the clean competition & \textbf{Yes} \\
\end{longtable}
}

The second is the \emph{premise} of the derivation, not its conclusion. Our measurements show that premise failing in a corpus-dependent way: on the controlled corpus \textbf{2 injected passages per stratum (0.1\% of 1,824)} already reach \texttt{poison@k} = 0.688, whereas on BBQ \textbf{356 per stratum (2\% of 17,792)} reach \textbf{0.000}. Nor is \(j^{*} \le k = 5\) sufficient on the controlled corpus, where 5 passages reach only 0.750 --- replacement requires out-ranking a large field of equally relevant clean passages. The defensible statement is therefore conditional: \textbf{once the attacker's passages are retrieved, violating an \(\varepsilon\)-constraint costs at most \(k\) replacements; whether they are retrieved is governed by the attack's relevance advantage, which is corpus- and encoder-dependent (Findings 4 and 7).} The error mattered: taken at face value it predicts that injection \emph{rate} is irrelevant, which our rate sweep falsifies.

\subsection{Proposition 3 --- the trilemma}\label{proposition-3-the-trilemma}

\textbf{Assumption (relevance overlap).} Every injected passage admits an admissible selection: for each \(a \in \mathcal{A}\) there is \(\mathcal{D} \in \mathcal{P}\) with \(a \in \mathcal{D}\), where \(\mathcal{P}\) is the set of selections meeting a utility floor (e.g.~every selected passage is within a relevance loss \(\delta\) of the unconstrained top-\(k\)). This is the standard assumption of corpus-poisoning work and our measurements support it --- the projection attack attains \texttt{poison@k} = 1.000 on our retrievers and 0.903--1.000 on real encoders at \(\rho = 0.1\%\).

\textbf{Proposition 3.} Under relevance overlap, for every \(\varepsilon\) there exists an injection of at most \(k\) passages per group that is simultaneously \(\varepsilon\)-feasible and \(\mathcal{P}\)-admissible.

\emph{Proof.} Prop. 1 supplies an injection that is \(\varepsilon\)-feasible for all \(\varepsilon\); relevance overlap supplies a \(\mathcal{P}\)-admissible selection containing each of its passages; the intersection is non-empty. \(\square\)

\textbf{Consequence --- the trilemma.} A distribution-constrained defense faces three requirements and can satisfy at most two:

{\def\LTcaptype{none} % do not increment counter
\begin{longtable}[]{@{}
  >{\raggedright\arraybackslash}p{(\linewidth - 2\tabcolsep) * \real{0.5000}}
  >{\raggedright\arraybackslash}p{(\linewidth - 2\tabcolsep) * \real{0.5000}}@{}}
\toprule\noalign{}
\begin{minipage}[b]{\linewidth}\raggedright
Requirement
\end{minipage} & \begin{minipage}[b]{\linewidth}\raggedright
Formal condition
\end{minipage} \\
\midrule\noalign{}

\bottomrule\noalign{}

\textbf{Soundness} --- exclude adversarial passages & \(\mathcal{R}_\varepsilon \cap \{\mathcal{D} : \mathcal{D}\cap\mathcal{A} \neq \emptyset\} = \emptyset\) \\
\textbf{Usefulness} --- admit good clean selections & \(\mathcal{R}_\varepsilon \cap \mathcal{P} \neq \emptyset\) \\
\textbf{Tightness} --- a non-vacuous constraint & \(\varepsilon < 1\) \\
\end{longtable}
}

Soundness fails for every \(\varepsilon\) whenever the attacker can construct an \(\varepsilon\)-feasible injection --- which requires only knowledge of \(p^{*}\) and \(\varepsilon\), both public for an auditable defense. So a defense must either (i) accept adversarial passages, (ii) restrict \(\mathcal{R}_\varepsilon\) so severely that \(\mathcal{R}_\varepsilon \cap \mathcal{P}\) shrinks or empties, sacrificing Usefulness, or (iii) set \(\varepsilon\) so large the constraint is vacuous, sacrificing Tightness.

\textbf{This is the formal content of the paper's central claim, and our experiments land on all three horns.} (i) Soundness fails: \texttt{poison@k} equals no defense at every \(\varepsilon\), on six retrievers, two corpora and all injection rates. (ii) Usefulness is squeezed: tightening the R1 budget to \(\varepsilon=0\) on the controlled corpus changes \texttt{in\_pool\_rate} from 0.017 to 0.088 with no security gain, and on the dense back-end from 0.2437 to 0.2812 --- i.e.~the constraint perturbs the selection without excluding anything. (iii) Tightness is vacuous at \(\varepsilon = 1\), where the constraint reduces to no defense by construction.

\subsection{\texorpdfstring{6.5 Proposition 4 --- why the family is \emph{structurally} exposed, not merely mis-tuned}{6.5 Proposition 4 --- why the family is structurally exposed, not merely mis-tuned}}\label{proposition-4-why-the-family-is-structurally-exposed-not-merely-mis-tuned}

Propositions 1--3 show that a distribution constraint cannot exclude the injection. A reader may reasonably ask whether that is a property of \emph{composition} constraints specifically, or of this defense family more broadly --- and in particular whether the non-compositional members, which the literature presents as the more robust alternatives, escape it. They do not, and the reason is worth stating separately because it covers the family as a whole.

\textbf{Proposition 4 (randomised defense under a known distribution).} Let a defense's selection mechanism be a fixed distribution \(\pi(\mathcal{D} \mid q, \mathcal{C})\) over retrieved sets, public to the attacker, and let \(U(\mathcal{D})\) be the attacker's utility (e.g.~the indicator that \(\mathcal{D}\) contains an injected passage). If any set in the support of \(\pi\) has \(U = 1\), the attacker can drive \(\mathbb{E}_{\pi}[U] \to 1\) by choosing an injection that places adversarial passages in every set with non-negligible probability mass.

\emph{Proof sketch.} Since \(\pi\) is fixed and known, \(\mathbb{E}_\pi[U]\) is computable in closed form as a function of the injection, and the attacker maximises it by search or by gradient. Randomisation does not bound \(\mathbb{E}_\pi[U]\); it bounds the \emph{variance} of \(U\) across realisations. An attacker whose success is measured in expectation over queries --- the usual case --- is indifferent to that variance. \(\square\)

The proposition is elementary, and that is the point: \textbf{randomisation defends against unpredictability, not against knowledge.} Two corollaries identify the cases that arise in this literature.

\textbf{Corollary 1 (public-distribution defenses).} Multi-query consistency {[}4{]} perturbs the query by token dropout at rate \(\delta\) and aggregates over \(n_v\) variants. The perturbation is random, but its \emph{rate} and its \emph{count} are public constants, so the expected aggregate score of a passage is a deterministic, differentiable function of the injected embedding which the attacker can maximise directly. Adding variants does not change this: the attacker optimises the expectation, and with many variants the empirical mean converges to it. The adaptive sweep reports the predicted collapse, from 0.729 at \(\lambda = 0\) to 1.000 by \(\lambda = 0.5\).

\textbf{Corollary 2 (penalty-observable defenses).} A defense that down-weights passages by a penalty \(f(d)\) computable from public information hands the attacker the constraint \texttt{}remain in the low-penalty region''. Whenever that constraint can be satisfied while raising similarity --- possible whenever the penalty is a continuous function of a perturbation the attacker controls --- the defense imposes a \emph{cost} rather than a \emph{barrier}. Off-manifold filtering is exactly this case, and the measured consequence is that it is the strongest defense statically (0.250 against 0.750 for no defense) and is defeated by \(\lambda = 0.25\).

We do not claim these observations are novel in general; they are the standard reason adaptive evaluation is required, and they are why security venues expect it. The claim is that \textbf{this defense family is structurally exposed to them}, because its two dominant mechanisms --- public randomisation and an observable penalty --- are precisely the two cases the proposition covers. Proposition 3 then explains why the remaining members fail for an independent reason, so the family is squeezed from both sides: composition constraints cannot exclude an individually admissible passage, and the non-compositional alternatives are defeated by an informed attacker. That is why we treat the failure as structural rather than as a matter of tuning, and why the honest remedy is detection (§6.6) rather than a better objective.

\subsection{Dimensionality does not help, and the only escape is individual-level evidence}\label{dimensionality-does-not-help-and-the-only-escape-is-individual-level-evidence}

Two consequences of the above are the actionable content of this section.

\textbf{(a) Adding dimensions does not shrink the admissible set usefully.} Extending the constraint from R1 to R2 to \(m\) dimensions does not remove adversarial selections from \(\mathcal{R}\); it gives the attacker more statistics to match, all of them computable from the public references. Prop. 1 applies verbatim with \(\mathcal{R}\) the \(m\)-dimensional feasible set. The attacker's cost rises, their \emph{feasibility} does not fall --- and since Prop. 2 bounds the cost by \(k\), the rise is cheap.

\textbf{(b) The only escape is evidence about individual passages.} To exclude an adversarial passage a defense must evaluate a property of \emph{that passage} that is not a function of the retrieved set's aggregate. In the adversarial setting this is a provenance question, answerable only from traces the attacker's optimisation leaves behind. §7 reports one candidate and its validity threat.

This reframes the problem. Prior work in this area is an \textbf{optimisation} problem --- which composition is fairest, and how to reach it. Under an adversary it becomes a \textbf{detection} problem, and §5.3 shows the two are not interchangeable: a defense can improve its fairness objective while its security objective is untouched. Prop. 3 is the reason the optimisation framing cannot be repaired by tuning: the failure is not at the optimum, it is in the feasible set.

\textbf{What this analysis does not establish.} Prop. 3 rests on the relevance-overlap assumption, which we support empirically but do not prove. A defense able to certify that an injected passage is \emph{not} \(\delta\)-competitive --- for instance by proving a relevance upper bound for it --- would escape the trilemma, and we know of no such defense in this literature. Constructing one is, in our view, the most promising direction this analysis opens. Prop. 2's bound also assumes replacement, i.e.~that injected passages out-rank those they displace; when they do not the bound is vacuous rather than wrong, which is exactly the corpus- and encoder-dependence documented in Findings 4 and 7. Finally, the analysis is pointwise in the query; a defense enforcing the constraint in expectation across queries has more room than the version analysed here, though ours are pointwise, so the measurements are unaffected.

\begin{center}\rule{0.5\linewidth}{0.5pt}\end{center}

\section{An Exploratory Provenance Signal}\label{an-exploratory-provenance-signal}

Proposition 1 indicates that a usable defense needs individual-level evidence. We report an exploratory candidate and are explicit about its validity status.

\textbf{Hypothesis.} A passage injected to be retrieved for many queries --- the projection attack optimises a \emph{single} direction \(\bar{q}\) so that one passage serves many queries --- will be retrieved for an anomalously large number of distinct queries. A legitimate passage is relevant to a small, specific set. The number of distinct queries retrieving a passage is computable at inference time from retrieval logs and requires no labels.

\textbf{Measurement.} We score each passage by the fraction of queries in the query set for which it appears in the top-\(k\).

At a threshold of the clean 90th percentile, detection rate is 1.000 with 3.9\% (BM25) and 0.2\% (dense) false positives.

\textbf{This result must not be taken at face value, and we state the reason.} The clean 90th percentile is \textbf{0.0000} --- 90\% of legitimate passages are retrieved by \emph{no} query. That is not a property of realistic corpora; it is a property of our construction, in which injected passages share a fixed query-aligned vocabulary that legitimate passages lack. The separation may therefore measure \emph{lexical overlap with our query template set} rather than \emph{retrieval optimisation}. In a natural corpus this signal may vanish entirely.

\textbf{Falsification protocol.} Before this can be claimed as a defense:
1. replicate on a natural corpus (TREC 2022 Fair Ranking Track {[}17{]}, Natural Questions) with injected passages generated by a query-agnostic attack {[}3{]} so that no shared vocabulary is introduced by construction;
2. report the signal's AUC with confidence intervals, and the resulting detection/false-positive trade-off;
3. verify that threshold selection does not require knowledge of the injection rate.

We report the signal because it is the direction Proposition 1 points to, and because a negative result from step (1) would itself be informative about which provenance cues survive in realistic settings.

\begin{center}\rule{0.5\linewidth}{0.5pt}\end{center}

\section{Discussion and Limitations}\label{discussion-and-limitations}

\textbf{What the failure is, and what it is not.} Our results do not show that existing fairness defenses are implemented incorrectly, nor that their reported gains are illusory in the setting they target. They show something narrower and, we think, more useful: the statistic each defense optimises is insensitive to the adversary that motivates the defense. Concretely, {[}5{]}'s optimisation, {[}6{]}'s proportion adjustment, {[}7{]}'s embedder control, and {[}8{]}'s exposure equalisation all act on aggregate composition. Under natural corpus skew --- the setting those works study --- this is appropriate, and their improvements are real. Under an adversary, §5.2 shows the same machinery leaves adversarial inclusion unchanged, and §6 shows why no refinement along that axis can change it. The distinction matters for how the field reads this result: the claim is about \textbf{evaluation validity}, and it implies a change in what gets measured rather than a rewrite of what has been built.

\textbf{An error we made twice, and why that is the point.} §3.2 identifies a class of mistakes that is easy to make and hard to notice: R2 statistics evaluated against a per-query clean retrieval, or against a corpus reference, are dominated by topic-conditioned stance variation that the attack preserves, and therefore report a healthy system while the retrieved set is fully adversarial. We made that error, corrected it by adopting the cross-group gap, and then made a second error of the same family one level up in §5.7: the cross-group gap is itself an absolute difference, and it is blind to an attack that relocates both groups rather than separating them. Both errors share a signature that makes them dangerous --- they return a plausible, healthy-looking number rather than an obviously wrong one. We therefore make two recommendations rather than one. Any stance-style metric should be specified together with its reference \emph{and} accompanied by a control showing it responds to a known attack; and where a metric is computed as a difference between two groups, the per-group values should be reported alongside it, because the difference can be flat while both components move a long way.

\textbf{Adversarial evaluation is a measurement problem before it is a defense problem.} Taken together, the failures in this paper --- a composition statistic that counts what the attack balances, a reference-based statistic that measures topic rather than stance, an absolute difference that cancels a joint shift, a self-report attribution probe that a model answers affirmatively by construction, an entailment convention whose reversed form returns zeros, and a null reported without an equivalence bound --- are all failures of instrument rather than of method. None of them would be caught by improving an attack or a defense. We suspect this generalises beyond retrieval fairness: any evaluation whose statistic is an aggregate over the quantity an adversary controls is vulnerable to the same class of error. The mitigation that worked here was mechanical rather than conceptual --- a control that flips the signal and checks the statistic moves, an equivalence bound on every null, and a per-component report for every difference.

\textbf{The three failure modes, and what each one predicts about work we did not run.} §3.4 stated the taxonomy from the form of the statistics; §5 confirmed it on measurements. Collecting the confirmations makes the framework's reach explicit, including for defenses we did not implement.

\emph{F1, balanced count.} Confirmed twice: the R1 statistics are invariant under pairwise injection (§5.1), and the positive control shows the same statistics detect a 3000-passage composition change with the corpus size held fixed (Table 1). So the invariance is a property of the attack, not the instrument. \textbf{The prediction for work we did not run}: any defense whose objective is a group count or share --- the re-ranking of {[}5{]}, the proportion adjustment of {[}6{]}, the embedder control of {[}7{]}, the exposure equalisation of {[}8{]} --- inherits this invariance exactly, without regard to how well it is optimised, because the injection is constructed to satisfy whatever composition the defense prefers. Our measurements on re-implementations of that objective support this and we have no reason to expect a different outcome from a better-optimised version of the same statistic.

\emph{F2, preserved nuisance.} Confirmed twice, and the second confirmation is the stronger one: \texttt{stance\_div} and \texttt{stance\_onesided} are flat across the clean and fully attacked conditions (§5.1), \emph{and} they are flat across the positive control where the true composition swings from 0.022 to 0.978 (Table 1). A statistic that does not move when the ground truth moves will not move for an adversary either. \textbf{The prediction}: any metric anchored to a per-query clean retrieval or a corpus-level rate is measuring the query's topic conditioning, which the attack preserves by construction, so evaluations that report such a metric alongside a clean baseline are reporting a quantity orthogonal to the threat.

\emph{F3, cancelled shift.} Confirmed at two layers, and this is the one we got wrong ourselves: the cross-group gap is blind to the retrieval-layer attack (§5.1) and to the generation-layer attack (§5.7), and on the generation layer our first analysis used it anyway and reported a null. \textbf{The prediction}: any difference, gap or variance across groups is invariant under a common-mode move, so a defense or an evaluation that reports only a gap will under-report any attack that relocates groups together rather than separating them --- which is what a matched-pair injection does by construction.

The three modes also compose, which is why the defense family fails as a family rather than in parts. F1 accounts for the composition-constrained members; §6.4 proves that no refinement along that axis escapes it. F3 accounts for the non-compositional members that report a group difference. And §6.5 shows the remaining mechanisms --- public randomisation and an observable penalty --- fail for a fourth reason that is not a measurement failure at all: an informed attacker optimises the expectation of a public distribution, and treats a computable penalty as a constraint rather than a barrier.

\textbf{The constructive half.} A taxonomy of failure modes is only useful if something escapes it, and the three modes share one structural feature: each is an invariance of an \emph{aggregation} over groups. Reporting the per-group quantities \(\phi_g\) and their paired changes is not invariant under any of the three, because it makes no aggregation claim for the adversary to preserve. That is a cheap change to how a result is reported, it requires no new defense, and §5.7 is an existence proof that it works: on identical retrieval conditions, the aggregated gap moved by 0.002 while the per-group levels moved by 0.13 at \(p < 0.01\). The recommendation is therefore not \texttt{}build a better defense metric'' but \texttt{}stop reporting only the aggregate''.

\textbf{A caveat we insist on, and a worked example of it from this paper.} The per-group report is more sensitive than an aggregate, and sensitivity is not validity. The clearest example is our own: the free-form probe reports a large, replicated per-group shift on the controlled corpus and nothing on the naturally written one, and it would have been easy to publish that as a finding about corpora. It is not --- the probe scores answers that take no position in 99.0\% of cases on natural text, so it was structurally incapable of reporting the effect that a commitment-forcing probe then found there. An instrument that manufactures commitment will report an effect where there is none; an instrument that suppresses it will report a null that does not exist. We hit the second case, and the control that caught it was measuring commitment rather than trusting the margin. We therefore report both probes and both corpora, and treat the disagreement between them as the result.

\textbf{Limitations.}

\begin{enumerate}
\def\labelenumi{\arabic{enumi}.}
\item
  \textbf{The measurement requires group- and stance-annotated passages, and natural retrieval corpora do not have them.} This is the binding constraint on this line of work, and it is worth stating plainly rather than as a data-collection to-do. R1 needs a group label per passage; R2 needs a stance label per passage. NQ and MS MARCO --- the standard neutral retrieval corpora --- carry neither, and stance toward a social group is not a property that can be inferred from a neutral passage, because a neutral passage does not express one. We therefore could not substitute a neutral corpus for BBQ, and neither can the defenses we evaluate: an R1-constraining defense needs the same group labels, and an R2-constraining defense needs labels that no existing neutral corpus provides. The framework's applicability is thus bounded by annotation, not by computation, and constructing a neutral, group- and stance-annotated retrieval benchmark is the field's outstanding infrastructure problem. Our two corpora bound the behaviour from the two sides available today: a template-controlled corpus with an exactly known reference, and a naturally written bias benchmark. We flag this as an open problem rather than a gap in the present study.
\item
  \textbf{Neither corpus is a neutral retrieval benchmark, and they disagree in two respects.} The text-only template attack is effective only on the controlled corpus (Finding 4), and the R2 constraint reduces adversarial inclusion only on BBQ. We report both conditions rather than selecting the favourable one. The mechanism behind the first disagreement is identified (lexical room depends on corpus repetitiveness), which we regard as a finding rather than a gap; the second is a smaller effect whose magnitude we expect to depend on the corpus's pre-existing stance imbalance.
\item
  \textbf{The absolute R2 metric does not transfer.} \texttt{stance\_gap} is exactly 0 on the balanced controlled corpus and near its maximum (0.9326) on BBQ \emph{before any attack}, where it is dominated by the benchmark's own stereotyped construction. On such a corpus only the change from the clean baseline is informative. A corpus-relative normalisation of \(\Delta_{\text{R2}}\) would be preferable and we leave it open.
\item
  \textbf{Injection budget is a rate, not a count.} We report an injection rate throughout (§5.4), because a fixed passage count measures corpus size and produced a spurious conclusion in our own earlier experiment. Readers comparing against work that reports absolute counts should convert.
\item
  \textbf{Retrieval back-ends, and what the scale check changed.} Nine retrievers are evaluated: BM25, a self-contained feature-hashing dense retriever, SPLADE and SPLADE-large, and five dense semantic encoders --- GTE-base, GTE-large, Contriever, E5-base-v2 and E5-large-v2 (§5.5, §5.6). An earlier draft of this paper listed the large checkpoints as \texttt{}to be added'' and \emph{predicted} that they would disagree with the base models. We have since measured it and report the correction: the prediction was right in direction and far too weak in magnitude, since \textbf{GTE-base \(\rightarrow\) GTE-large moves from 0.0625 to 0.5000 --- an eight-fold increase in susceptibility within a single encoder family and training recipe.} That has three consequences for how this study should be read. First, feature-hashing dense retrieval retains the lexical attack surface and should not be read as a stand-in for a semantic encoder. Second, text-attack effectiveness varies \textbf{nine-fold across the base dense encoders alone} (Contriever 0.5625 versus GTE-base and E5-base-v2 at 0.0625) and the susceptible one is the encoder used by the fair-ranking work we compare against {[}8{]}. Third, and most importantly, \textbf{even holding the family and training recipe fixed, changing the checkpoint changes the answer qualitatively} --- so a single-encoder robustness claim does not establish how robust a defense is in general. This is a limitation of our study and, we argue, of the standard evaluation protocol in this area. We have not swept Contriever at large scale, nor multilingual or instruction-tuned embedders, and we expect the spread to widen rather than narrow in those directions. Our scale pairs are parameter-matched (110M \(\rightarrow\) 335M for GTE and E5; 66M \(\rightarrow\) 110M for SPLADE), so the SPLADE pair tests a smaller size delta than the other two and a \texttt{}large SPLADE is no more robust'' conclusion carries correspondingly less weight.
\item
  \textbf{The generation-stage evaluation is multi-generator and multi-probe, and the probes disagree in a way that is itself a result.} §5.7 reports four generators spanning three model families (API-served Qwen and DeepSeek, plus open-weight Mistral-7B run locally in 4-bit on an RTX 5060), two corpora, and two probes: free-form answers scored by entailment, and a forced-choice probe scored by the choice it forces. Three limits remain. (i) \emph{Retrieval conditions are measured on one backbone} (GTE-base), held fixed across the generation-stage comparison, so the encoder-dependence of §5.5 is not varied here. (ii) \emph{The two probes support different conclusions, and the freer one is the weaker.} Free-form answers commit to a position in only 15.6\% of cases on the controlled corpus and 1.0\% on BBQ, so that probe is near-blind on natural text; the forced-choice probe detects the effect on both. We report both rather than the favourable one, because a probe that manufactures commitment can also manufacture an effect, and §5.7.1 records the checks we ran to establish that the forced choice tracks statement content rather than option position (100\% versus 35\% agreement with the placement of the favourable statement). (iii) \emph{The effect is generator-dependent.} One of three generators carries most of the signal on both corpora, a second shows a single group on BBQ, and the third shows none reaching significance there. A result of this shape is not summarisable as one number, and we no longer report one. Whether a larger generator would show more or less is an open question: we ran one open-weight model and three API models, and the scale of the API models is not controlled.

  \emph{Query diversity.} The controlled corpus's 48 retrieval slots correspond to only \textbf{16 distinct question texts} (each repeated 3\(\times\), once per paraphrase variant). Any statistic aggregated over the 48 is therefore aggregated over 16 questions, and the effective sample size is a third of what the table suggests. The natural-corpus evaluation does not have this problem (BBQ supplies 144 distinct, naturally written questions), which is an additional reason to prefer it for output-level measurement rather than treating it as a replication check.

  \emph{The generator does not follow an inverted context on this corpus.} We ran a negative control that flips the stance of every retrieved passage while holding the question fixed. On the controlled corpus the answer was \textbf{unchanged in 6 of 6 cases under inversion}, but \textbf{changed in 6 of 6 cases when the context was replaced by unrelated passages} or removed entirely. The generator is therefore reading the context at the level of \emph{topic} --- it knows a question about an engineering evaluation calls for naming an engineer --- but not at the level of \emph{stance}, and the answers are determined by the question's framing rather than by the evidence on trial. This is a property of a template corpus whose questions presuppose their own answer, and it means a null propagation result on that corpus would have been uninformative. The forced-choice probe reported in §5.7 avoids the problem by construction, which is why it is the primary measurement; we now report the control because the free-form probe alone cannot distinguish \texttt{}no propagation'' from \texttt{}generator ignores stance''.

  \emph{A metric we reported as failed is now fixed, and the fix is not cosmetic.} The per-passage attribution probe asked the generator itself whether its answer relied on a given passage; the judge answered YES almost always and the resulting EAE-D was a constant 1.000 in every condition. The cause is direction: an NLI model must be asked whether the \textbf{passage entails the answer}, not the reverse. Calibrated on this corpus, a passage the answer used scores 0.996 in the correct direction and 0.008 in the reversed one, and an unused passage scores 0.000 in both --- so the reversed form returns noise rather than an obviously wrong value. §5.7's EAE-D figures should be read as produced by the probe we have now replaced.

  A multi-generator evaluation with entailment-based attribution is in progress and is not reported here; we do not otherwise report attributed exposure {[}8{]} or generator bias {[}4, 7{]}.
\item
  \textbf{Binary groups.} Following {[}6, 7, 8{]}, we use two groups per stratum. Extension to \(|\mathcal{G}| > 2\) is mechanical for R1 and R2 but is not evaluated here. We note that the race/ethnicity category of BBQ is markedly imbalanced in our corpus build (960 passages about the protected group versus 88 about the non-protected group), which is itself a property of the benchmark worth flagging for anyone reusing it.
\end{enumerate}

\begin{center}\rule{0.5\linewidth}{0.5pt}\end{center}

\section{Conclusion}\label{conclusion}

Fairness defenses for retrieval-augmented generation are chosen and validated by aggregate statistics of the retrieved set. We have shown that the standard statistics are structurally insensitive to the attack they are meant to detect, that the insensitivity has one identifiable cause, and that the defense failure which follows is a consequence of the measurement failure rather than a separate phenomenon.

The cause is that the attack balances what these statistics count and relocates what they difference. Pairwise poisoning --- the natural adversarial construction, in which injected passages reuse the legitimate template inventory in matched pairs --- is balanced in group composition by construction, so every composition statistic reads clean. And because it moves both groups rather than opening a gap between them, any statistic computed as an absolute difference between groups is dominated by the corpus's pre-existing asymmetry and is likewise blind. The diagnostic is the per-group shift. We identify this at the retrieval layer, where two of three candidate stance statistics are flat across the clean and fully attacked conditions, and then again at the generation layer in our own first measurement of propagation, where the absolute gap moved by less than 0.002 while the per-group shifts moved by 0.13 at \(p \le 0.0035\). Reporting the second occurrence against ourselves is part of the argument: the error is not exotic, it is the natural first thing to compute.

The practical consequence is that the R1-only defense is not weak but inert, and we state this as a ceiling rather than as a null that further data might overturn. The unconstrained baseline already admits adversarial passages on every query where it can, so no reduction is available even in principle; the per-query difference is identically zero and the equivalence bound is exactly ±0.0000. That is not because the constraint cannot change the selection --- it can, and does --- but because changing which passages are selected does not displace a single adversarial one. Extending the constraint to a second dimension improves the statistic it measures without changing adversarial inclusion. The problem is therefore not one of optimisation but of detection, and it requires evidence about individual passages rather than about the set they belong to.

We close with the protocol the results imply. The first four points are about \textbf{measurement validity} --- whether a statistic can see the attack at all --- and the last five about \textbf{adversarial evaluation hygiene}, which is what makes a robustness claim interpretable:

\begin{enumerate}
\def\labelenumi{\arabic{enumi}.}
\tightlist
\item
  \textbf{Measure per-group shifts, not cross-group differences.} An absolute gap between groups is dominated by pre-existing corpus asymmetry and is blind to an attack that relocates both groups together. Report both group-level quantities, and report the change in each.
\item
  \textbf{Report injection as a rate, not a count.} A fixed passage count measures corpus size rather than attack strength, and produced a spurious conclusion in our own earlier experiment.
\item
  \textbf{Report the encoder as a factor.} Susceptibility is a per-checkpoint property that is not monotone in encoder size --- within one family and training recipe, scaling GTE raised text-attack success eight-fold while scaling E5 lowered it to zero. A single-encoder robustness claim reports an unmeasured property of that checkpoint.
\item
  \textbf{Give equivalence bounds for null claims.} \texttt{}No effect detected'' and \texttt{}no effect'' are different claims, and only the second supports a conclusion about a defense family. Where the per-query difference is identically zero this is easy and conclusive; where it is not, the bound is the honest statement.
\item
  \textbf{State the threat model explicitly, including whether the attacker knows the defense.} A robustness claim without this qualifier is uninterpretable, and the qualifier is what separates the two evaluations in §5.2.
\item
  \textbf{Report adversarial inclusion as a function of attacker strength, not at a single operating point.} One value cannot distinguish a robust defense from one evaluated below its failure threshold. In our sweep every defense fails by \(\lambda \le 2\); a sweep stopping at \(\lambda = 0.25\) would have shown three apparently robust defenses.
\item
  \textbf{Report a utility-preserving baseline.} A defense can achieve low adversarial inclusion trivially by returning fewer or worse passages, which is why §5.2 reports \texttt{in\_pool\_rate} alongside \texttt{poison@k}.
\item
  \textbf{For a randomised defense, state the distribution and assume it is known.} Report the attacker's optimal response to the expectation rather than to a sample; by Proposition 4, randomisation bounds variance and not the expectation.
\item
  \textbf{For a penalty-based defense, report whether the penalty is computable from public information.} If it is, treat it as a constraint in the attacker's optimisation rather than as an unknown; by Corollary 2 that is the difference between a cost and a barrier.
\end{enumerate}

Point 9 is the one most easily overlooked and, in our experiments, decisive for the defense that looked strongest statically. Point 5 is the one most often omitted in this literature. We suggest that points 1--4 are prerequisites for a fairness statistic to be reported at all under adversarial conditions, and points 5--9 for a robustness claim to be credited.

We have reported an exploratory provenance signal in the direction the protocol implies --- evidence about individual passages rather than about aggregates --- together with the reason it may not survive replication on a natural corpus. Establishing which provenance cues do survive is, in our view, the central open problem this work exposes.

\begin{center}\rule{0.5\linewidth}{0.5pt}\end{center}

\section{Data and Code Availability}\label{data-and-code-availability}

The implementation, configuration files, and the scripts that regenerate every number in §5 are released at \textbf{https://github.com/luj00133-dev/rag-poison-fairness}. The controlled corpus is generated deterministically from a fixed seed and requires no dataset download; the BBQ corpus is reconstructed from the official benchmark files by a released loader that derives stance labels from BBQ's own annotations. §5.7 additionally requires a DeepSeek API key, supplied through the \texttt{DEEPSEEK\_API\_KEY} environment variable and never stored in the repository.

\textbf{How to regenerate.}

\begin{verbatim}
# retrieval layer (CPU only, no API key)
python -m src.run_experiment --config configs/default.json      # ~25 s
python -m src.run_experiment --config configs/bbq.json          # ~180 s
python -m src.run_experiment --config configs/align_gte_ad.json # GTE-base, ~12 min
python -m src.run_experiment --config configs/align_e5.json     # E5-base, ~10 min
python -m src.run_experiment --config configs/contriever_ad.json # Contriever, ~3 min
python -m src.run_experiment --config configs/align_splade.json # SPLADE distil, ~4 min

# encoder-scale check of §5.6 (large checkpoints, CPU only)
python -m src.run_experiment --config configs/align_gte_large.json      # ~23 min
python -m src.run_experiment --config configs/align_e5_large.json       # ~10 min
python -m src.run_experiment --config configs/align_splade_large.json   # ~12 min

# the analyses behind §5.5 and §5.6
python analysis/compare_six_backbones.py
python analysis/compare_scale.py
python analysis/diagnose_gte_large.py     # embedding geometry
python analysis/diagnose_poison_scores.py # the mechanism behind Table C1

# generation layer (requires DEEPSEEK_API_KEY)
python -m src.run_attribution --config configs/attribution.json
\end{verbatim}

Real-encoder runs require \texttt{HF\_ENDPOINT=https://hf-mirror.com} on networks where huggingface.co is unreachable.

\textbf{API key.} The key used for the §5.7 runs is rotated and revoked before publication; the repository contains no credentials, and \texttt{results/attribution\_cache.json} (a prompt\(\rightarrow\)reply cache) is regenerated locally rather than shipped.

\begin{center}\rule{0.5\linewidth}{0.5pt}\end{center}

\section{Appendix A. Reproduction}\label{appendix-a.-reproduction}

\begin{verbatim}
# controlled-corpus suite (CPU only; ~25 s)
python -m src.run_experiment --config configs/default.json

# natural-corpus replication on BBQ (CPU only; ~165 s)
python -m src.run_experiment --config configs/bbq.json

# development-scale run (~5 s)
python -m src.run_experiment --config configs/default.json --quick

# corpus comparison table
python compare_corpora.py

# over-generalisation probe
python src/probe_overgeneralisation.py --quick --retriever bm25
\end{verbatim}

The BBQ data is downloaded from the official repository:

\begin{verbatim}
python -m src.data.bbq_loader          # builds and reports corpus statistics
\end{verbatim}

(Expects \texttt{data/bbq/\{Gender\_identity,Disability\_status,Age,Race\_ethnicity\}.jsonl}
from \texttt{https://github.com/nyu-mll/BBQ/tree/main/data}.)

Outputs: \texttt{results/\textless{}run\_tag\textgreater{}/\{per\_query,aggregate,by\_stratum,adaptive\}.csv} and \texttt{summary.txt}.

The full retrieval-layer study runs on CPU: the controlled suite completes in
25 s and the BBQ replication in 165 s on a desktop CPU, with no GPU required.
The backbone alignment of §5.5 and §5.6 adds roughly 75 minutes of CPU time
across eight encoder checkpoints, the slowest single run being GTE-large at
23 minutes; the largest model loaded is a 335M-parameter encoder, which needs
under 2 GB in fp32. No GPU is used anywhere in this paper.

\begin{center}\rule{0.5\linewidth}{0.5pt}\end{center}

\section{Appendix B. Paired tests and confidence intervals}\label{appendix-b.-paired-tests-and-confidence-intervals}

Every number in §5 is an aggregate over queries. Because the paper's central claim is a null, we give the paired tests on the underlying per-query data here, including the equivalence bound for each null cell. Reproduce with the script \texttt{analysis/stats\_paper.py}.

\textbf{Table B1.} Paired tests of R1-constraint inertness against no defense. Controlled corpus, \texttt{template\_plus\_projection}, \(\rho = 0.5\%\), per-query adversarial-passage inclusion.

{\def\LTcaptype{none} % do not increment counter
\begin{longtable}[]{@{}
  >{\raggedright\arraybackslash}p{(\linewidth - 14\tabcolsep) * \real{0.1250}}
  >{\raggedright\arraybackslash}p{(\linewidth - 14\tabcolsep) * \real{0.1250}}
  >{\raggedright\arraybackslash}p{(\linewidth - 14\tabcolsep) * \real{0.1250}}
  >{\raggedright\arraybackslash}p{(\linewidth - 14\tabcolsep) * \real{0.1250}}
  >{\raggedright\arraybackslash}p{(\linewidth - 14\tabcolsep) * \real{0.1250}}
  >{\raggedright\arraybackslash}p{(\linewidth - 14\tabcolsep) * \real{0.1250}}
  >{\raggedright\arraybackslash}p{(\linewidth - 14\tabcolsep) * \real{0.1250}}
  >{\raggedright\arraybackslash}p{(\linewidth - 14\tabcolsep) * \real{0.1250}}@{}}
\toprule\noalign{}
\begin{minipage}[b]{\linewidth}\raggedright
Retriever
\end{minipage} & \begin{minipage}[b]{\linewidth}\raggedright
Defense
\end{minipage} & \begin{minipage}[b]{\linewidth}\raggedright
\(\varepsilon\)
\end{minipage} & \begin{minipage}[b]{\linewidth}\raggedright
n
\end{minipage} & \begin{minipage}[b]{\linewidth}\raggedright
\(\Delta\) mean
\end{minipage} & \begin{minipage}[b]{\linewidth}\raggedright
95\% CI
\end{minipage} & \begin{minipage}[b]{\linewidth}\raggedright
p
\end{minipage} & \begin{minipage}[b]{\linewidth}\raggedright
excluded effect
\end{minipage} \\
\midrule\noalign{}

\bottomrule\noalign{}

bm25 & repr\_group & 0.0 & 48 & +0.0000 & {[}+0.0000, +0.0000{]} & 1.000 & ±0.0000 \\
bm25 & repr\_group & 1.0 & 48 & +0.0000 & {[}+0.0000, +0.0000{]} & 1.000 & ±0.0000 \\
bm25 & repr\_both & 0.0 & 48 & +0.0000 & {[}+0.0000, +0.0000{]} & 1.000 & ±0.0000 \\
bm25 & repr\_both & 1.0 & 48 & +0.0000 & {[}+0.0000, +0.0000{]} & 1.000 & ±0.0000 \\
bm25 & repr\_stance & 0.0 & 48 & +0.0000 & {[}+0.0000, +0.0000{]} & 1.000 & ±0.0000 \\
bm25 & repr\_stance & 1.0 & 48 & +0.0000 & {[}+0.0000, +0.0000{]} & 1.000 & ±0.0000 \\
dense & repr\_group & 0.0 & 48 & +0.0000 & {[}+0.0000, +0.0000{]} & 1.000 & ±0.0000 \\
dense & repr\_group & 1.0 & 48 & +0.0000 & {[}+0.0000, +0.0000{]} & 1.000 & ±0.0000 \\
dense & repr\_both & 0.0 & 48 & +0.0000 & {[}+0.0000, +0.0000{]} & 1.000 & ±0.0000 \\
dense & repr\_both & 1.0 & 48 & +0.0000 & {[}+0.0000, +0.0000{]} & 1.000 & ±0.0000 \\
dense & repr\_stance & 0.0 & 48 & +0.0000 & {[}+0.0000, +0.0000{]} & 1.000 & ±0.0000 \\
dense & repr\_stance & 1.0 & 48 & +0.0000 & {[}+0.0000, +0.0000{]} & 1.000 & ±0.0000 \\
\end{longtable}
}

Where the constraint changes nothing, the per-query difference vector is \textbf{identically zero}, not merely small: the bootstrap CI collapses to {[}0, 0{]} and the permutation test is degenerate by construction. The informative statistic is therefore the equivalence bound --- the smallest effect the data exclude --- and because the differences are exactly zero rather than approximately zero, these cells exclude \emph{any} effect on adversarial-passage inclusion, not merely one above a threshold. The \texttt{repr\_stance} rows are the contrast case: there the constraint does change the selection, the difference vector is non-degenerate, and both the CI and the p-value are meaningful.

The point of Table B1 is the last column. For the R1-only and joint constraints the difference is zero per query, so the CI is {[}0, 0{]} and the equivalence bound is ±0.0000: the data exclude any effect on adversarial inclusion, not merely an effect above some threshold. The \texttt{repr\_stance} rows are included as a contrast, since there the constraint genuinely changes the selection.

\begin{center}\rule{0.5\linewidth}{0.5pt}\end{center}

\section{Appendix C. Secondary tables}\label{appendix-c.-secondary-tables}

Tables moved out of the main body to keep it readable. Each is referenced from the section that discusses it.

\textbf{Table C1.} Adaptive attacker (defense-aware), \texttt{poison@k}. Controlled corpus, \(\rho = 2\%\).

{\def\LTcaptype{none} % do not increment counter
\begin{longtable}[]{@{}
  >{\raggedright\arraybackslash}p{(\linewidth - 14\tabcolsep) * \real{0.1250}}
  >{\raggedright\arraybackslash}p{(\linewidth - 14\tabcolsep) * \real{0.1250}}
  >{\raggedright\arraybackslash}p{(\linewidth - 14\tabcolsep) * \real{0.1250}}
  >{\raggedright\arraybackslash}p{(\linewidth - 14\tabcolsep) * \real{0.1250}}
  >{\raggedright\arraybackslash}p{(\linewidth - 14\tabcolsep) * \real{0.1250}}
  >{\raggedright\arraybackslash}p{(\linewidth - 14\tabcolsep) * \real{0.1250}}
  >{\raggedright\arraybackslash}p{(\linewidth - 14\tabcolsep) * \real{0.1250}}
  >{\raggedright\arraybackslash}p{(\linewidth - 14\tabcolsep) * \real{0.1250}}@{}}
\toprule\noalign{}
\begin{minipage}[b]{\linewidth}\raggedright
Defense
\end{minipage} & \begin{minipage}[b]{\linewidth}\raggedright
\(\lambda\)=0
\end{minipage} & \begin{minipage}[b]{\linewidth}\raggedright
\(\lambda\)=0.25
\end{minipage} & \begin{minipage}[b]{\linewidth}\raggedright
\(\lambda\)=0.5
\end{minipage} & \begin{minipage}[b]{\linewidth}\raggedright
\(\lambda\)=1
\end{minipage} & \begin{minipage}[b]{\linewidth}\raggedright
\(\lambda\)=2
\end{minipage} & \begin{minipage}[b]{\linewidth}\raggedright
\(\lambda\)=4
\end{minipage} & \begin{minipage}[b]{\linewidth}\raggedright
static rank
\end{minipage} \\
\midrule\noalign{}

\bottomrule\noalign{}

off-manifold filtering & \textbf{0.250} & \textbf{0.562} & \textbf{0.812} & 0.938 & 1.000 & 1.000 & 1st \\
\texttt{repr\_both} (R1+R2) & 0.625 & 0.875 & 1.000 & 1.000 & 1.000 & 1.000 & 2nd \\
multi-query consistency & 0.729 & 0.958 & 1.000 & 1.000 & 1.000 & 1.000 & 3rd \\
\end{longtable}
}

\textbf{Table C2.} The first generation-layer measurement, with the forced-choice probe on a single generator. Superseded by Table 12; kept because the contrast between the two is the point. Controlled corpus, GTE-base retrieval, \(\rho = 2\%\).

{\def\LTcaptype{none} % do not increment counter
\begin{longtable}[]{@{}
  >{\raggedright\arraybackslash}p{(\linewidth - 8\tabcolsep) * \real{0.2000}}
  >{\raggedright\arraybackslash}p{(\linewidth - 8\tabcolsep) * \real{0.2000}}
  >{\raggedright\arraybackslash}p{(\linewidth - 8\tabcolsep) * \real{0.2000}}
  >{\raggedright\arraybackslash}p{(\linewidth - 8\tabcolsep) * \real{0.2000}}
  >{\raggedright\arraybackslash}p{(\linewidth - 8\tabcolsep) * \real{0.2000}}@{}}
\toprule\noalign{}
\begin{minipage}[b]{\linewidth}\raggedright
Condition
\end{minipage} & \begin{minipage}[b]{\linewidth}\raggedright
fav. rate (g1)
\end{minipage} & \begin{minipage}[b]{\linewidth}\raggedright
fav. rate (g2)
\end{minipage} & \begin{minipage}[b]{\linewidth}\raggedright
\textbf{stance gap}
\end{minipage} & \begin{minipage}[b]{\linewidth}\raggedright
\(\Delta\) vs clean
\end{minipage} \\
\midrule\noalign{}

\bottomrule\noalign{}

clean & 0.562 & 0.708 & \textbf{0.2708} & --- \\
poisoned & 0.583 & 0.979 & \textbf{0.3958} & \textbf{+46\%} \\
\texttt{r1only} (R1) & 0.604 & 1.000 & \textbf{0.3958} & +46\% \\
\texttt{r2both} (R1+R2) & 0.583 & 1.000 & \textbf{0.4167} & +54\% \\
\end{longtable}
}

\textbf{Table C3.} Prop. 1 prediction vs.~measurement. R1-only constraint (\texttt{repr\_group}), strongest attack, largest injection rate per corpus; \(\Delta\) = constrained minus no defense.

{\def\LTcaptype{none} % do not increment counter
\begin{longtable}[]{@{}
  >{\raggedright\arraybackslash}p{(\linewidth - 10\tabcolsep) * \real{0.1667}}
  >{\raggedright\arraybackslash}p{(\linewidth - 10\tabcolsep) * \real{0.1667}}
  >{\raggedright\arraybackslash}p{(\linewidth - 10\tabcolsep) * \real{0.1667}}
  >{\raggedright\arraybackslash}p{(\linewidth - 10\tabcolsep) * \real{0.1667}}
  >{\raggedright\arraybackslash}p{(\linewidth - 10\tabcolsep) * \real{0.1667}}
  >{\raggedright\arraybackslash}p{(\linewidth - 10\tabcolsep) * \real{0.1667}}@{}}
\toprule\noalign{}
\begin{minipage}[b]{\linewidth}\raggedright
Corpus / retriever
\end{minipage} & \begin{minipage}[b]{\linewidth}\raggedright
\(\varepsilon\)=1.0
\end{minipage} & \begin{minipage}[b]{\linewidth}\raggedright
\(\varepsilon\)=0.5
\end{minipage} & \begin{minipage}[b]{\linewidth}\raggedright
\(\varepsilon\)=0.25
\end{minipage} & \begin{minipage}[b]{\linewidth}\raggedright
\(\varepsilon\)=0.1
\end{minipage} & \begin{minipage}[b]{\linewidth}\raggedright
\(\varepsilon\)=0.0
\end{minipage} \\
\midrule\noalign{}

\bottomrule\noalign{}

controlled / BM25 & 0.000 & 0.000 & 0.000 & 0.000 & 0.000 \\
controlled / feature-hash dense & 0.000 & 0.000 & 0.000 & 0.000 & 0.000 \\
BBQ / BM25 & 0.000 & 0.000 & 0.000 & 0.000 & 0.000 \\
BBQ / feature-hash dense & 0.000 & 0.000 & 0.000 & 0.000 & 0.000 \\
controlled / GTE-base & 0.000 & --- & --- & --- & 0.000 \\
controlled / Contriever & 0.000 & --- & --- & --- & 0.000 \\
\end{longtable}
}

\textbf{Table C4.} Over-generalisation signal.

{\def\LTcaptype{none} % do not increment counter
\begin{longtable}[]{@{}
  >{\raggedright\arraybackslash}p{(\linewidth - 10\tabcolsep) * \real{0.1667}}
  >{\raggedright\arraybackslash}p{(\linewidth - 10\tabcolsep) * \real{0.1667}}
  >{\raggedright\arraybackslash}p{(\linewidth - 10\tabcolsep) * \real{0.1667}}
  >{\raggedright\arraybackslash}p{(\linewidth - 10\tabcolsep) * \real{0.1667}}
  >{\raggedright\arraybackslash}p{(\linewidth - 10\tabcolsep) * \real{0.1667}}
  >{\raggedright\arraybackslash}p{(\linewidth - 10\tabcolsep) * \real{0.1667}}@{}}
\toprule\noalign{}
\begin{minipage}[b]{\linewidth}\raggedright
Retriever
\end{minipage} & \begin{minipage}[b]{\linewidth}\raggedright
Set
\end{minipage} & \begin{minipage}[b]{\linewidth}\raggedright
n
\end{minipage} & \begin{minipage}[b]{\linewidth}\raggedright
mean
\end{minipage} & \begin{minipage}[b]{\linewidth}\raggedright
median
\end{minipage} & \begin{minipage}[b]{\linewidth}\raggedright
p90
\end{minipage} \\
\midrule\noalign{}

\bottomrule\noalign{}

BM25 & poison & 16 & 0.1914 & 0.1875 & 0.2500 \\
BM25 & clean & 608 & 0.0032 & \textbf{0.0000} & \textbf{0.0000} \\
dense + projection (\(\lambda\)=1) & poison & 16 & 0.3086 & 0.2188 & 0.6250 \\
dense + projection (\(\lambda\)=1) & clean & 608 & 0.0001 & \textbf{0.0000} & \textbf{0.0000} \\
\end{longtable}
}

\textbf{Table C5.} Bootstrap 95\% CIs on the positive claims (per-question values, 10,000 resamples, no defense).

{\def\LTcaptype{none} % do not increment counter
\begin{longtable}[]{@{}
  >{\raggedright\arraybackslash}p{(\linewidth - 8\tabcolsep) * \real{0.2000}}
  >{\raggedright\arraybackslash}p{(\linewidth - 8\tabcolsep) * \real{0.2000}}
  >{\raggedright\arraybackslash}p{(\linewidth - 8\tabcolsep) * \real{0.2000}}
  >{\raggedright\arraybackslash}p{(\linewidth - 8\tabcolsep) * \real{0.2000}}
  >{\raggedright\arraybackslash}p{(\linewidth - 8\tabcolsep) * \real{0.2000}}@{}}
\toprule\noalign{}
\begin{minipage}[b]{\linewidth}\raggedright
Condition
\end{minipage} & \begin{minipage}[b]{\linewidth}\raggedright
Retriever
\end{minipage} & \begin{minipage}[b]{\linewidth}\raggedright
Metric
\end{minipage} & \begin{minipage}[b]{\linewidth}\raggedright
mean
\end{minipage} & \begin{minipage}[b]{\linewidth}\raggedright
95\% CI
\end{minipage} \\
\midrule\noalign{}

\bottomrule\noalign{}

clean & bm25 & R2 stance gap & 0.0000 & {[}0.0000, 0.0000{]} \\
clean & bm25 & R1 drift (TV) & 0.0000 & {[}0.0000, 0.0000{]} \\
clean & dense & R2 stance gap & 0.0000 & {[}0.0000, 0.0000{]} \\
clean & dense & R1 drift (TV) & 0.0000 & {[}0.0000, 0.0000{]} \\
template & bm25 & R2 stance gap & 0.7436 & {[}0.6068, 0.8718{]} \\
template & bm25 & R1 drift (TV) & 0.1625 & {[}0.1062, 0.2208{]} \\
template & dense & R2 stance gap & 0.6167 & {[}0.4556, 0.7667{]} \\
template & dense & R1 drift (TV) & 0.0812 & {[}0.0438, 0.1208{]} \\
template\_plus\_projection & bm25 & R2 stance gap & 0.7436 & {[}0.6068, 0.8718{]} \\
template\_plus\_projection & bm25 & R1 drift (TV) & 0.1625 & {[}0.1062, 0.2208{]} \\
template\_plus\_projection & dense & R2 stance gap & 1.0000 & {[}1.0000, 1.0000{]} \\
template\_plus\_projection & dense & R1 drift (TV) & 0.1625 & {[}0.1187, 0.2083{]} \\
\end{longtable}
}

\begin{center}\rule{0.5\linewidth}{0.5pt}\end{center}

\section{References}\label{references}
\begin{thebibliography}{20}
\bibitem{ref1} P. Lewis, E. Perez, A. Piktus, et al.~Retrieval-augmented generation for knowledge-intensive NLP tasks. \emph{NeurIPS}, 2020.
\bibitem{ref2} W. Zou, R. Geng, B. Wang, J. Jia. PoisonedRAG: Knowledge corruption attacks to retrieval-augmented generation of large language models. \emph{USENIX Security}, 2025.
\bibitem{ref3} B. Zhang, Y. Chen, Z. Liu, et al.~Practical poisoning attacks against retrieval-augmented generation. \emph{arXiv:2504.03957}, 2025.
\bibitem{ref4} L. Wang, T. Zhu, L. Qin, L. Gao, W. Zhou. Bias amplification in RAG: Poisoning knowledge retrieval to steer LLMs. \emph{IEEE Transactions on Dependable and Secure Computing}, 23(5):11033--11050, 2026. (arXiv:2506.11415)
\bibitem{ref5} Y. Zhao, V. Efthymiou, J. Nummenmaa, K. Stefanidis. Fairness-aware retrieval optimization for retrieval-augmented generation. \emph{arXiv:2605.15790}, 2026.
\bibitem{ref6} X. Wu, S. Li, H.-T. Wu, Z. Tao, Y. Fang. Does RAG introduce unfairness in LLMs? Evaluating fairness in retrieval-augmented generation systems. \emph{COLING}, pp.~10021--10036, 2025. (arXiv:2409.19804)
\bibitem{ref7} T. Kim, J. M. Springer, A. Raghunathan, M. Sap. Mitigating bias in RAG: Controlling the embedder. \emph{Findings of ACL}, 2025. (arXiv:2502.17390)
\bibitem{ref8} T. E. Kim, F. Diaz. Towards fair RAG: On the impact of fair ranking in retrieval-augmented generation. \emph{ICTIR}, pp.~33--43, 2025. (arXiv:2409.11598)
\bibitem{ref9} J. Liang, Y. Wang, C. Li, et al.~GraphRAG under fire. \emph{IEEE Symposium on Security and Privacy}, 2026. (arXiv:2501.14050)
\bibitem{ref10} H. Ha, Q. Zhan, J. Kim, et al.~MM-PoisonRAG: Disrupting multimodal RAG with local and global poisoning attacks. \emph{arXiv:2502.17832}, 2025.
\bibitem{ref11} Y. Liu, Z. Yuan, G. Tie, et al.~Poisoned-MRAG: Knowledge poisoning attacks to multimodal retrieval augmented generation. \emph{arXiv:2503.06254}, 2025.
\bibitem{ref12} K. Edemacu, M. M. Shokri. Hidden in the metadata: Stealth poisoning attacks on multimodal retrieval-augmented generation. \emph{arXiv:2603.00172}, 2026.
\bibitem{ref13} G. Bagwe, S. S. Chaturvedi, X. Ma, et al.~Your RAG is unfair: Exposing fairness vulnerabilities in retrieval-augmented generation via backdoor attacks. \emph{arXiv:2509.22486}, 2025.
\bibitem{ref14} S. K. Mohanty, R. Patel, K. Yuvaraj, et al.~TriShieldRAG: 3 rings, one blind spot in layered defenses for retrieval-augmented generation. \emph{arXiv:2607.23838}, 2026.
\bibitem{ref15} A. Parrish, A. Chen, N. Nangia, et al.~BBQ: A hand-built bias benchmark for question answering. \emph{Findings of ACL}, pp.~2086--2105, 2022.
\bibitem{ref16} M. Nadeem, A. Bethke, S. Reddy. StereoSet: Measuring stereotypical bias in pretrained language models. \emph{ACL}, 2021.
\bibitem{ref17} M. D. Ekstrand, G. McDonald, A. Raj, I. Johnson. Overview of the TREC 2022 fair ranking track. \emph{TREC}, 2022.
\bibitem{ref18} S. Dai, X. Chen, S. Xu, L. Pang, Z. Dong, J. Xu. Bias and unfairness in information retrieval systems: New challenges in the LLM era. \emph{KDD}, pp.~6437--6447, 2024.
\bibitem{ref19} M. Hu, H. Wu, Z. Guan, et al.~No free lunch: Retrieval-augmented generation undermines fairness in LLMs, even for vigilant users. \emph{arXiv:2410.07589}, 2024.
\bibitem{ref20} R. Shrestha, Y. Zou, Q. Chen, et al.~FairRAG: Fair human generation via fair retrieval augmentation. \emph{CVPR}, 2024.
\end{thebibliography}
\end{document}
